\documentclass[10pt,a4paper]{article}
\usepackage{amsmath,amssymb,amsthm,graphicx}
\usepackage[a4paper,landscape,margin=18mm]{geometry}
\usepackage{mathtools,amssymb,amsthm,enumitem,tikz-cd,etoolbox,thmtools}
\usepackage{unicode-math}
\setmathfont{latinmodern-math.otf}
\usepackage[capitalize]{cleveref}
\newcounter{chapter}
\NewDocumentCommand{\chapter}{sm}{\IfBooleanTF{#1}{\section*{#2}}{\refstepcounter{chapter}\section{#2}}}
\let\oiint\relax
\newcommand{\co}{\!:}
\newcommand{\define}[1]{\index{#1}\emph{#1}} % définition du mot #1
\newcommand{\subdefine}[2]{\index{#1!#2}\emph{#2}} % définition du mot #2 comme sous-entrée de #1
\newcommand{\definex}[2]{\index{#2}\emph{#1}} % affiche #1 index #2
\newcommand{\subdefinex}[3]{\index{#2!#3}\emph{#1}} % affiche #1 index #3 comme sous-entrée de #2

\newcommand{\image}[1]{{\centering\includegraphics{images/#1}\par}}

\newcommand{\A}{\mathcal{A}} % un atlas
\newcommand{\E}{\mathcal{E}} % un espace affine
\newcommand{\Cinf}{C^\infty}
\newcommand{\Rel}{\mathcal{R}} % une relation différentielle

\newcommand{\rst}[2]{#1|_{#2}}
\newcommand{\bigO}[1]{O\left(#1\right)}

\DeclareMathOperator{\Diff}{Diff}
\DeclareMathOperator{\Loop}{\mathcal{L}}
\DeclareMathOperator{\Bij}{Bij}
\DeclareMathOperator{\Aut}{Aut}
\DeclareMathOperator{\Id}{Id}
\DeclareMathOperator{\rk}{rk}
\DeclareMathOperator{\GL}{GL}
\DeclareMathOperator{\SO}{SO}
\DeclareMathOperator{\SU}{SU}
\DeclareMathOperator{\Lin}{L}
\DeclareMathOperator{\Hom}{Hom}
\DeclareMathOperator{\Bsym}{B_\mathrm{s}}
\DeclareMathOperator{\Hess}{Hess}
\DeclareMathOperator{\im}{im}
\DeclareMathOperator{\coker}{coker}
\DeclareMathOperator{\supp}{supp} % support
\DeclareMathOperator{\ind}{ind} % indice d'un point critique
\DeclareMathOperator{\sgn}{ε} % signature d'une permutation
\DeclareMathOperator{\Alt}{Alt} % antisymétrisation
\DeclareMathOperator{\Div}{div} % divergence
\DeclareMathOperator{\tr}{tr} % trace
\DeclareMathOperator{\pr}{pr} % une projection
\DeclareMathOperator{\Int}{Int} % intérieur
\DeclareMathOperator{\Conv}{Conv} % enveloppe convexe
\DeclareMathOperator{\IntConv}{IntConv} % intérieur enveloppe convexe
\DeclareMathOperator{\Conn}{Conn} % connected component
\DeclareMathOperator{\Span}{Span} % sev engendré
\DeclareMathOperator{\Gr}{Gr} % Grassmanienne
\DeclareMathOperator{\CP}{CP} % Corrugation process
\DeclareMathOperator{\Sec}{Sec} % Sections
\DeclareMathOperator{\Hol}{Hol} % Holonomic sections
\DeclareMathOperator{\codim}{codim} % codimension
\DeclareMathOperator{\set}{\mathcal{P}} % powerset

% Cohomologie de de Rham, version avec indice dR ou non
%\newcommand{\HdR}[1][*]{H^{#1}_{\!\scriptscriptstyle d\!R}}
%\newcommand{\HdRc}[1][*]{H^{#1}_{\!\scriptscriptstyle d\!R, c}}
\newcommand{\HdR}[1][*]{H^{#1}}
\newcommand{\HdRc}[1][*]{H^{#1}_c}

\newcommand{\transp}[1]{#1^T}

\newcommand{\eE}{\mathcal{E}}
\newcommand{\fF}{\mathcal{F}}
\newcommand{\tT}{\mathcal{T}}

\newcommand{\X}{\mathfrak{X}} % champs de vecteurs

\newcommand{\C}{C} % cube
\newcommand{\F}{\mathcal{F}} % non-holonomic section
\newcommand{\HH}{\mathcal{H}} % homotopy of non-holonomic section
\newcommand{\G}{\mathcal{G}} % non-holonomic section

\newcommand{\pt}{\mathrm{pt}} % point

\DeclareMathOperator{\Op}{Op} % non-specified open neighborhood
\DeclareMathOperator{\dCzero}{d_{C^0}} % C^0 distance
\DeclareMathOperator{\dCone}{d_{C^1}} % C^1 distance
\DeclareMathOperator{\bs}{bs} % base map of a jet bundle section

% Linear maps between vector spaces
\renewcommand{\L}[2]{L(#1, #2)}
% Update a linear map #2 using dual pair #1 and vector #3
\newcommand{\Upd}[3]{\Upsilon_{#1}\left(#2, #3\right)}
\newcommand{\IT}[1]{\mathcal{I}_{#1}}
% Déclaration de théorème avec thmtools
\declaretheorem[numberwithin=chapter]{theorem}
\declaretheorem[sibling=theorem]{proposition}
\declaretheorem[sibling=theorem]{corollary}
\declaretheorem[sibling=theorem]{remark}
\declaretheorem[sibling=theorem]{lemma}
\declaretheorem[sibling=theorem]{definition}
\declaretheorem[sibling=theorem]{example}

\declaretheorem[numbered=no, name=Theorem]{theorem-intro}
\declaretheorem[numbered=no, name=Proposition]{proposition-intro}
\declaretheorem[numbered=no, name=Corollary]{corollary-intro}
\declaretheorem[numbered=no, name=Lemma]{lemma-intro}
\declaretheorem[numbered=no, name=Definition]{definition-intro}

% Cache les commandes plastex
\newcommand{\uses}[1]{}
\newcommand{\proves}[1]{}
\newcommand{\lean}[1]{}
\newcommand{\leanok}{}

%\newcommand{\D}[1]{\gls{D}[#1]}
\newcommand{\D}[1]{\mathrm{Diff}(#1)}

\newenvironment{diagram}[1][]{\[\begin{tikzcd}[#1]}{\end{tikzcd}\]}
\newenvironment{dessin}[1][]{\[\begin{tikzpicture}[#1]}{\end{tikzpicture}\]}

% Bugfix :
% http://tex.stackexchange.com/questions/262617/issue-with-xrightarrow
\makeatletter
  \patchcmd{\arrowfill@}{-7mu}{-14mu}{}{}
  \patchcmd{\arrowfill@}{-7mu}{-14mu}{}{}
  \patchcmd{\arrowfill@}{-2mu}{-4mu}{}{}
  \patchcmd{\arrowfill@}{-2mu}{-4mu}{}{}
\makeatother

\newcommand{\oiint}{\iint}

\definecolor{fond}{RGB}{255,255,255}

\newcommand{\N}[1]{\mathcal{N}_{\!\!#1}} % neighborhood
\title{Full relative parametric C0-dense h-principle for open ample first-order relations}
\author{Patrick Massot, Floris van Doorn, Oliver Nash}
\date{github:leanprover-community-sphere-eversion:df3e74a3fbe1ae42c2cd9b075d77ff3863c1b795; Apache-2.0}
\begin{document}
\maketitle
\noindent\textit{Editorial note.} Original author mathematical source is retained without mathematical corrections or generated proof text. Chapter presentation is adapted to the article wrapper; Blueprint Lean/proves/uses metadata is hidden by the original print macros.\par\medskip
\noindent\textit{Editorial note.} The selected counted target is the full open-ample h-principle. Sphere eversion and the alternate local argument are preserved as authored applications, not additional counts of the same core proof.\par\medskip
\noindent\textit{Editorial note.} SOURCE\_RESOLUTION confirms images/\#1 is an unused image macro, with no concrete image dependencies; no source assets have been invented.\par\medskip
\noindent\textit{Allowed standard prerequisites (editorial).}\begin{itemize}
\item Finite-dimensional smooth calculus, partitions of unity, elementary convex geometry and smooth manifold/jet bundle definitions are cited as ordinary prerequisites exactly as in the author Blueprint.
\end{itemize}\medskip
\chapter*{Introduction}
\label{cha:intro}

This project had two goals. First we wanted to check whether a proof
assistant can do differential topology. Many people still think that
formal mathematics are mostly suitable for algebra, combinatorics, or
foundational studies. So we chose one of the most famous examples
of geometric topology theorems associated to tricky geometric intuition:
the existence of sphere eversions. Note however that we won't focus on
any of the many videos of explicit sphere eversions. We will prove a
general theorem which immediately implies the existence of sphere
eversions.

The second goal of this project was to experiment using a formalization
blueprint that evolves with the project until we get a proof that has
very closely related formal and informal presentations. A full proof (by
normal pen and paper standards) was written before the formalization
effort began. This proof evolved a lot during the formalization. In
particular, the chapter on the global theory required a lot of work
during the formalization in order to ensure that its technical lemmas
are both fully correct and actually sufficient for our purposes.

In this introduction, we will describe the mathematical context of this
project, the main definitions and statements, and outline the proof
strategy.

Gromov observed that it's often fruitful to distinguish two kinds of geometric
construction problems. He says that a geometric construction problem satisfies
the $h$-principle if the only obstructions to the existence of a solution come
from algebraic topology. In this case, the construction is called flexible,
otherwise it is called rigid. This definition is purposely vague. We
will see a rather general way to give it a precise meaning, but one must keep
in mind that such a precise meaning will fail to encompass a number of
situations that can be illuminated by the $h$-principle dichotomy point of view.

The easiest example of a flexible construction problem which is not totally
trivial and is algebraically obstructed is the deformation of immersions
of circles into planes. Let $f_0$ and $f_1$ be two maps from $𝕊^1$ to $ℝ^2$
that are immersions. Since $𝕊^1$ has dimension one, this mean that both
derivatives $f_0'$ and $f_1'$ are nowhere vanishing maps from $𝕊^1$ to
$ℝ^2$. The geometric object we want to construct is a (smooth) homotopy
of immersions from $f_0$ to $f_1$, ie a smooth map $F \co 𝕊^1 × [0, 1] → ℝ^2$
such that $\rst{F}{𝕊^1 × \{0\}} = f_0$, $\rst{F}{𝕊^1 × \{1\}} = f_1$, and
each $f_p := \rst{F}{𝕊^1 × \{p\}}$ is an immersion. If such a homotopy
exists then, $(t, p) ↦ f_p'(t)$ is a homotopy from $f_0'$ to $f_1'$ among
maps from $𝕊^1$ to $ℝ^2 ∖ \{0\}$. Such maps have a well defined winding
number $w(f'_i) ∈ ℤ$ around the origin, the degree of the normalized map
$f'_i/‖f'_i‖ \co 𝕊^1 → 𝕊^1$. So $w(f_0') = w(f_1')$ is a necessary condition
for the existence of $F$, which comes from algebraic topology.  The
Whitney--Graustein theorem states that this necessary condition is also
sufficient. Hence this geometric construction problem is flexible.  One
can give a direct proof of this result, but it also
follows from general results proved in this project (although we
haven't formalized this consequence of our work).

An important lesson from the above example is that algebraic topology
can give us more than a necessary condition. Indeed the (one-dimensional)
Hopf degree theorem ensures that, provided $w(f'_0) = w(f'_1)$, there
exists a homotopy $g_p$ of nowhere vanishing maps relating $f'_0$ and
$f'_1$. We also know from the topology of $ℝ^2$ that $f_0$ and $f_1$ are
homotopic, say using the straight-line homotopy
$p ↦ f_p = (1-p)f_0 + pf_0$. But there is no a priori relation
between $g_p$ and the derivative of $f_p$ for $p ∉ \{0, 1\}$.
So we can restate the crucial part of the
Whitney--Graustein theorem as: there is a homotopy of immersion from
$f_0$ to $f_1$ as soon as there is (a homotopy from $f_0$ to $f_1$) and
a homotopy from $f_0'$ to $f_1'$ among nowhere vanishing maps. The
parenthesis in the previous sentence indicated that this condition is
always satisfied, but it is important to keep in mind for
generalizations. Gromov says that such a homotopy of uncoupled pairs
$(f, g)$ is a formal solution of the original problem.

One can generalize this discussion of uncoupled maps replacing a
map and its derivative for maps from a manifold $M$
to a manifold $N$. The so called $1$-jet space $J^1(M, N)$ is the
space of triples $(m, n, φ)$ with $m ∈ M$, $n ∈ N$, and
$φ ∈ \Hom(T_mM, T_nN)$, the space of linear maps from $T_mM$ to $T_nN$.
One can define a smooth manifold structure on $J^1(M, N)$, of dimension
$\dim(M) + \dim(N) + \dim(M)\dim(N)$ which fibers over $M$, $N$ and
their product $J^0(M, N) := M × N$. Beware that the notation $(m, n, φ)$
does not mean that $J^1(M, N)$ is a product of three manifolds, the space
where $φ$ lives depends on $m$ and $n$.
Any smooth map $f \co M → N$ gives rise to a section $j^1f$ of
$J^1(M, N) → M$ defined by $j^1f(m) = (m, f(m), T_m f)$. Such a section is
called a \emph{holonomic section} of $J^1(M, N)$.
In the Whitney--Graustein example, we use the canonical trivialization of
$T𝕊^1$ and $Tℝ^2$ to represent $j^1f$ has a pair of maps $(f, f')$. The
role played by $(f, g)$ in this example is played in general by sections
of $J^1(M, N) → M$ which are not necessarily holonomic.

One can generalize this discussion to $J^r(M, N)$ which remembers
derivatives of maps up to order $r$ for some given $r ≥ 0$. One can also
consider sections of an arbitrary bundle $E → M$ instead of functions
from $M$ to $N$, which are sections of the trivial bundle $M × N → N$.
But the case of $J^1(M, N)$ is sufficient for our purposes.

\begin{definition-intro}
  A first order differential relation $\Rel$ for maps from $M$ to
  $N$ is a subset of $J^1(M, N)$. A solution of $\Rel$ is a function
  $f \co M \to N$ such that $j^1f(m)$ is in $\Rel$ for all $m$.
  A \emph{formal solution} of $\Rel$ is a non-necessarily holonomic
  section of $J^1(M, N) → M$ which takes value in $\Rel$.

  The partial differential relation $\Rel$ satisfies the $h$-principle
  if any formal solution $σ$ of $\Rel$ is homotopic, among formal
  solutions, to some holonomic one $j^1f$.
\end{definition-intro}

For instance, an immersion of $M$ into $N$ is a solution of
\[
\Rel = \{(m, n, φ) ∈ J^1(M, N) \;|\; φ \text{ is injective}\}.
\]
As we saw with
the Whitney--Graustein problem, we are not only interested to individual
solutions, but also in families of solutions. In differential
topology, a smooth family of maps between manifolds $X$ and $Y$ is a smooth map
$h \co P × X → Y$ seen as the collection of maps $h_p \co x ↦ h(p, x)$. Here $P$
stands for ``parameter space''. A smooth family of sections of $E → X$ is a
smooth family of maps $σ \co P × X → E$ such that each $σ_p$ is a section.

In such a case it is important that we start with a family of formal solutions
that is holonomic for some values of the parameter and we don't modify it for
those parameters. In the curve example $P = [0, 1]$, the formal solution is
holonomic for parameters $0$ and $1$, and we want to keep the start and end
curves. In this work we don't use manifolds with boundary when it is not
necessary so we rather use $ℝ$ as a parameter space.

More generally it can also happen that a family of formal solutions
$σ : P × M → J^1(M, N)$ has the property that $σ_p$ is holonomic at $m ∈ M$ for
some values of $p$ and $m$ and we want to preserve $σ$ near the corresponding
set in $P × M$. This leads to the following definition.


\begin{definition-intro}
  A partial differential relation $\Rel ⊂ J^1(M, N)$ satisfies the relative and
  parametric $h$-principle if every family of formal solutions
  $σ \co M × P → J^1(M, N)$ which are holonomic for $(p, m)$ near some closed
  set $C ⊂ P × M$, is homotopic to a family of holonomic sections, and this
  homotopy can be chosen constant near $C$.
\end{definition-intro}

One can also insist on the deformed solution to be $C^0$-close to the original
one. In this case one talks about a $C^0$-dense $h$-principle.

Using this vocabulary, we can state the Smale-Hirsch immersion theorem as saying
that the relation of immersions satisfies all forms of the $h$-principle provided
the dimension of the target manifold is larger than the dimension of the source.

This theorem covers the Whitney--Graustein theorem (in its second
form, assuming the existence of a homotopy between derivatives). But there are
much less intuitive applications. The most famous one is the existence of
sphere eversions: one can ``turn $𝕊^2$ inside-out among immersions of $𝕊^2$ into
$ℝ^3$).

\begin{corollary-intro}[Smale 1958]
	There is a homotopy of immersion of $𝕊^2$ into $ℝ^3$ from the inclusion map to
	the antipodal map $a \co q ↦ -q$.
\end{corollary-intro}

The reason why this is turning the sphere inside-out is that $a$ extends as
a map from $ℝ^3 ∖ \{0\} → ℝ^3 ∖ \{0\}$ by
\[
	\hat{a} \co q ↦ -\frac{1}{‖q‖^2}q
\]
which exchanges the interior and exterior of $𝕊^2$. More abstractly, one can say
the normal bundle of $𝕊^2$ is trivial, hence one can extend $a$ to a tubular
neighborhood of $𝕊^2$ as an orientation preserving map. Since $a$ is orientation
reversing, any such extension will be reversing co-orientation.

\begin{proof}[Proof of the sphere eversion corollary]
	We denote by $ι$ the inclusion of $𝕊^2$ into $ℝ^3$.
	We set $j_t = (1-t)ι	+ ta$. This is a homotopy from $ι$ to $a$ (but not an
	immersion for $t=1/2$).
  We need to check there is no obstruction to building a homotopy of formal
  solutions above those maps. One could show that the relevant homotopy group
  (replacing $π_1(𝕊^1)$ from the Whitney--Graustein example) is
  $\pi_2(\SO_3(ℝ))$. This group is trivial, hence there is no obstruction. But
  actually we can write an explicit homotopy here, without computing
  $\pi_2(\SO_3(ℝ))$.  Using the canonical trivialization of the tangent
	bundle of $ℝ^3$, we can set, for $(q, v) ∈ T𝕊^2$,
	$G_t(q, v) = \mathrm{Rot}_{Oq}^{πt}(v)$, the rotation around axis $Oq$ with
	angle $πt$. The family $σ \co t ↦ (j_t, G_t)$ is a homotopy of formal immersions
	relating $j^1ι$ to $j^1a$. The above theorem ensures this family is homotopic,
	relative to $t = 0$ and $t = 1$, to a family of holonomic formal immersions,
	ie a family $t ↦ j^1f_t$ with $f_0 = ι$, $f_1 = a$, and each $f_t$ is an
	immersion.
\end{proof}

The Smale-Hirsch theorem and its above corollary follows from a more general theorem:
the $h$-principle for open and ample first order differential relations (see
below). We will prove this theorem using a technique which is even more
general: convex integration. For instance this technique also underlies the
constructions of paradoxical isometric embeddings, which could be a nice
follow-up project.

We'll end this introduction by describing the key construction of convex
integration, since it is very nice and elementary. Convex integration
was invented by Gromov around 1970, inspired in particular by the
$C^1$ isometric embedding work of Nash and the original proof of
flexibility of immersions. This term is pretty vague however, and there
are several different implementations. The newest one, and by far the
most efficient one, is Mélanie Theillière's corrugation process from
2018. And this is what we will use.

Let $f$ be a map from $ℝ^n$ to $ℝ^m$. Say we want to turn $f$ into a
solution of some partial differential relation. For instance if we are
interested in immersions, we want to make sure its differential is everywhere
injective. We will ensure this by tackling each partial derivative in turn. In
the immersion example, we first make sure $∂_1f(x) := ∂f(x)/∂x_1$ is non-zero
for all $x$. Then we make sure $∂_2f(x)$ is not colinear to $∂_1f(x)$. Then we
make sure $∂_3f(x)$ is not in the plane spanned by the two previous
derivatives, etc\dots\ until all $n$ partial derivatives are everywhere
linearly independent.

In general, what happens is that, for each number $j$ between $1$ and
$n$, we wish $∂_jf(x)$ could live in some open subset
$Ω_x ⊂ ℝ^m$. Assume there is a smooth family of loops $γ \co ℝ^n × 𝕊^1 → ℝ^m$ such
that each $γ_x$ takes values in $Ω_x$, and has average value
$\int_{𝕊^1} γ_x = ∂_j f(x)$. Obviously such loops
can exist only if $∂_jf(x)$ is in the convex hull of $Ω_x$,
and we will see this condition is
almost sufficient. In the immersion case, this convex hull condition
will always be met because, from the above description, we see that
$Ω_x$ will always be the complement of a linear subspace with
codimension at least two.

For some large positive $N$, we replace $f$ by the new map
\[
  x ↦ f(x) + \frac1N  ∫_0^{Nx_j} \left[γ_x(s) - ∂_jf(x)\right]ds.
\]

A wonderfully easy exercise shows that, provided $N$ is large enough,
we have achieved $∂_j f(x) ∈ Ω_x$, almost without modifying
derivatives $∂_i f(x)$ for $i ≠ j$, and almost without moving
$f(x)$. See \ref{prop:theilliere} for a precise statement.
This technique is called convex integration
since we are taking an integral under the assumption that
$∂_jf(x)$ is in the convex hull of $Ω_x$.

In addition, if we assume that $γ_x$ is constant (necessarily with value
$∂_j f(x)$) for $x$ near some subset $K$ where
$∂_j f(x)$ was already good, then nothing changed on $K$ since the
integrand vanishes there. It is also easy to damp out this modification
by multiplying the integral by a cut-off function. So this is a very
local construction, and it isn't obvious how the absence of homotopical
obstruction, embodied by the existence of a formal solution, should
enter the discussion. The answer is that is essentially provides a way
to coherently choose base points for the $γ_x$ loops.

Now that we've seen how convex hulls enter the discussion we can provide one
last definition and state the actual main theorem that we formalized.

\begin{definition-intro}
  A relation $\Rel ⊂ J^1(M, N)$ is ample if, for every $(x, y, φ) ∈ \Rel$ and
  every hyperplane $H ⊂ T_xM$, the convex hull of the connected component of $φ$ in
  \[
    \left\{ψ ∈ \Hom(T_xM, T_yN) \;|\; \rst{ψ}{H} = \rst{φ}{H} \text{ and } (x, y, ψ) ∈ \Rel\right\}
  \]
  is the whole set of $ψ$ such that $\rst{ψ}{H} = \rst{φ}{H}$.
\end{definition-intro}

We can now state our goal in its full glory.

\begin{theorem-intro}[Gromov]
  For any manifolds $M$ and $N$, any relation $\Rel ⊂ J^1(M, N)$ that is open
  and ample satisfies the full $h$-principle (relative, parametric and $C^0$-dense).
\end{theorem-intro}

\Cref{chap:loops} provides the loops supply. \Cref{chap:local} then
discusses the local theory, including the key construction above,
and \cref{chap:global} finally moves to manifolds, and proves the main
theorem and its sphere eversion corollary.
\Cref{chap:local_eversion} explains how the first two chapters are already
enough to derive Smale's theorem, although in a slightly less natural way than
using the manifold case. This served as an intermediate target in the
formalization, and can be used for elementary teaching since it does not
require any theory of manifolds.

\chapter{Loops}
\label{chap:loops}

\section{Introduction}
\label{sec:loops_introduction}

In this chapter, we explain how to construct families of loops to feed into the
corrugation process explained at the end of the introduction.

Throughout this document, $E$ and $F$ will denote finite-dimensional
real vector spaces.

\begin{definition}
  \label{def:loop}
  \lean{Loop, Loop.average, Loop.support}
  \leanok
  A loop is a map defined on the circle $𝕊^1 = ℝ/ℤ$ with values in a
  finite-dimensional vector space.
  It can also freely be seen as $1$-periodic maps defined on $ℝ$.

  The average of a loop $γ$ is $\bar γ := \int_{𝕊^1} γ(s)\, ds$.

  The \emph{support} of a family $γ$ of loops in $F$ parametrized by $E$ is
  the closure of the set of $x$ in $E$ such that $γ_x$ is not a constant loop.
\end{definition}

All of this chapter is devoted to proving the following proposition.

\begin{proposition}
  \label{prop:∃_loops}
  \lean{exist_loops}
  \uses{def:loop}
  \leanok
  Let $K$ a compact set in $E$. Let $Ω$ be an open set in $E × F$.

  Let $β$ and $g$ be smooth maps from $E$ to $F$.
  Write $Ω_x := \{ y ∈ F \mid (x, y) ∈ Ω\}$, assume that $β(x) ∈ Ω_x$ for all $x$,
  and that $g(x) = β(x)$ near $K$.

  If, for every $x$, $g(x)$ is in
  the convex hull of the connected component of $Ω_x$ containing $β(x)$,
  then there exists a smooth family of loops
  \[
    γ \co E × [0, 1] × 𝕊^1 → F, (x, t, s) ↦ γ^t_x(s)
  \]
  such that, for all $x ∈ E$, all $t ∈ ℝ$ and all  $s ∈ 𝕊^1$,
  \begin{itemize}
    \item
      $γ^t_x(s) ∈ Ω_x$
    \item
      $γ^0_x(s) = γ^t_x(1) = β(x)$
    \item
      $\bar γ^1_x = g(x)$
    \item
      $γ^t_x(s) = β(x)$ if $x$ is near $K$.
  \end{itemize}
\end{proposition}

Let us briefly sketch the geometric idea behind the above proposition
if we pretend there is only one point $x$, and drop it from the
notation, and also focus only on $γ^1$.
By assumption, there is a finite collection of points $p_i$ in $Ω$ and $λ_i ∈
[0, 1]$ such that $g$ is the barycenter $\sum λ_i p_i$. Since $Ω$ is open and
connected, there is a smooth loop $γ_0$ which goes through each $p_i$. The
claim is that $g$ is the average value of $γ = γ_0 ∘ h$ for some
self-diffeomorphism $h$ of $𝕊^1$. The idea is to choose $h$ such that
$γ$ rushes to $p_1$, stays there during a time roughly $λ_1$, rushes to
$p_2$, etc. But, in order to achieve average exactly $g$, it seems like $h$
needs to be a discontinuous piecewise constant map. The assumption that $Ω$ is
open means that the convex hull is open, which gives enough slack to get away with
a smooth $h$.

In the previous proof sketch, there is a lot of freedom in constructing $γ$,
which is problematic when trying to do it consistently when $x$ varies.

\section{Surrounding points}
\label{sec:preliminaries}

This section collects elementary results about convex sets in finite
dimensional real vector spaces that will help to construct families of loops.
In this section, $E$ is a real vector space with (finite) dimension $d$.
The discussion will center around the following definition which is tailored to
our ulterior needs.

\begin{definition}
  \label{def:surrounds_points}
  \lean{SurroundingPts}
  \leanok
  A point $x$ in $E$
  is surrounded by points $p_0$, \dots, $p_d$ if those points are
  affinely independent and there exist weights $w_i ∈ (0, 1)$ with sum $1$
  such that $x = \sum_i w_i p_i$.
\end{definition}

Note that, in the above definition, the number of points $p_i$ is fixed
by the dimension $d$ of $E$, and that the weights $w_i$ are the barycentric
coordinates of $x$ with respect to the affine basis $p_0, \ldots, p_d$.

The first important point in this definition is that surrounding is
smoothly locally stable: if $x$ is surrounded by a collection of points $p$
then points that are close to $y$ are surrounded by every collection of points
$q$ that is closed to $p$, and the relevant barycentric coordinates smoothly
depend on $y$ and $q$. The precise statement follows.

\begin{lemma}
  \label{lem:smooth_barycentric_coord}
  \uses{def:surrounds_points}
  \lean{smooth_surrounding}
  \leanok
  For every $x$ in $E$ and every collection of points $p ∈ E^{d+1}$
  surrounding $x$, there is a function $w \co E × E^{d+1} → ℝ^{d+1}$ such that,
  for every $(y, q)$ in a neighborhood of $(x, p)$,
  \begin{itemize}
    \item
      $w$ is smooth at $(y, q)$
    \item
      $w(y, q) > 0$
    \item
      $\sum_{i=0}^d w_i(y, q) = 1$
    \item
      $y = \sum_{i=0}^d w_i(y, q)q_i$.
  \end{itemize}
\end{lemma}

\begin{proof}
  \leanok
  Let:
  \begin{align*}
    A = E \times \{ q \in E^{d+1} ~|~ \mbox{$q$ is an affine basis for $E$} \},
  \end{align*}
  and define:
  \begin{align*}
    w \co A &\to ℝ^{d+1}\\
    (y, q) &\mapsto \mbox{barycentric coordinates of $y$ with respect to $q$}.
  \end{align*}

  If we fix an affine basis $b$ of $E$, we may express $w$ as a
  ratio of determinants in terms of coordinates relative to $b$. More precisely,
  by Cramer's rule, if $0 \le i \le d$ and $w_i$ is the $i^{\rm th}$ component of $w$,
  then:
  \begin{align*}
    w_i (y, q) = \det M_i (y, q) / \det N (q)
  \end{align*}
  where $N(q)$ is the $(d+1)\times (d+1)$ matrix whose columns are the barycentric
  coordinates of the components of $q$ relative to $b$, and $M_i (y, q)$ is $N(q)$
  except with column $i$ replaced by the barycentric coordinates of $y$ relative
  to $b$.

  Since determinants are smooth functions and $(y, q) \mapsto \det N(q)$ is
  non-vanishing on $A$, $w$ is smooth on $A$.

  Finally define:
  \begin{align*}
    U = w^{-1}((0, \infty)^{d+1}),
  \end{align*}
  and note that $U$ is open in $A$, since it is the preimage of an open set under the
  continuous map $w$. In fact since $A$ is open, $U$ is open as a subset of $E \times E^{d+1}$.
  Note that $(x, p) \in U$ since $p$ surrounds $x$.

  We may extend $w$ to $E \times E^{d+1}$ by giving it any values at all outside $A$.
\end{proof}

Then we need a criterion ensuring a point $x$ is surrounded by a collection of points taken
in a given subset $P$. The first temptation is to hope that $x$ being in the interior of
the convex hull of $P$ is enough. But this is not true. For instance the center of a square
in a plane is in the interior of the convex hull of the set $P$ of vertices of the square,
but it isn't surrounded by any set of vertices. This counter example also shows that
the stability lemma above is slightly less trivial than it sounds.

The rest of this section is devoted to the following result that proves no such issue arises
when $P$ is open.

\begin{proposition}
  \label{prop:surrounded_by_open}
  \uses{def:surrounds_points}
  \lean{surrounded_of_convexHull}\leanok
  If a point $x$ of $E$ lies in the convex hull of an open set $P$,
  then it is surrounded by some collection of points belonging to $P$.
\end{proposition}

This proposition will be proven at the end of this section.
We'll first need the Carathéodory lemma:

\begin{lemma}[Carathéodory's lemma]
\label{lem:caratheodory}
\lean{convexHull_eq_union}
\leanok
  If a point $x$ of $E$ lies in the convex hull of a set $P$, then $x$
  belongs to the convex hull of a finite set of affinely independent points
  of $P$.
\end{lemma}

\begin{proof}
  \leanok
  By assumption, there is a finite set of points $t_i$ in $P$ and
  weights $f_i$ such that $x = \sum f_i t_i$, each $f_i$ is non-negative
  and $\sum f_i = 1$.
  Choose such a set of points of minimum cardinality. We argue by
  contradiction that such a set must be affinely independent.

  Thus suppose that there is some vanishing combination $\sum g_i t_i$ with
  $\sum g_i = 0$ and not all $g_i$ vanish.
  Let $S = \{i | g_i > 0\}$.
  Let $i_0$ in $S$ be an index minimizing $f_i/g_i$. We shall obtain our
  contradiction by showing that $x$ belongs to the convex hull of the set
  $\{t_i| i \ne i_0\}$, which has cardinality strictly smaller than
  $\{t_i\}$.

  We thus define new weights $k_i = f_i - g_i f_{i_0}/g_{i_0}$.
  These weights sum to $\sum f_i - (\sum g_i)f_{i_0}/g_{i_0} = 1$
  and $k_{i_0} = 0$.
  Each $k_i$ is non-negative, thanks to the choice of $i_0$ if
  $i$ is in $S$ or using that $f_i$, $-g_i$ and $f_{i_0}/g_{i_0}$
  are all non-negative when $i$ is not in $S$.
  It remain to compute
  \begin{align*}
    \sum_{i ≠ i_0} k_i t_i &= \sum_i k_i t_i \\
      &= \sum_i (f_i - g_i f_{i_0}/g_{i_0}) t_i  \\
      &= \sum_i f_i t_i - \left(\sum_i g_i t_i\right)f_{i_0}/g_{i_0})   \\
      &= x
  \end{align*}
  where we use $k_{i_0} = 0$ in the first equality.
\end{proof}



\begin{lemma}
  \label{lem:interior_chab}
  \lean{AffineBasis.interior_convexHull}\leanok
  Given an affine basis $b$ of $E$, the interior of the convex hull of
  $b$ is the set of points with strictly positive barycentric coordinates.
\end{lemma}

\begin{proof}
  \leanok
  For each $i$, let:
  \[
    w_i \co E \to ℝ
  \]
  be the $i^{\rm th}$ barycentric coordinate with respect to the basis $b$.
  Since $E$ is finite-dimensional, each $w_i$ is a continuous open map. For
  such a map, the operation of taking interior commutes with preimage, and so
  we have:
  \begin{align*}
    \IntConv(b) &= \Int\left(\bigcap_i w_i^{-1}([0, \infty))\right)\\
                &= \bigcap_i \Int(w_i^{-1}([0, \infty))\\
                &= \bigcap_i w_i^{-1}(\Int([0, \infty))\\
                &= \bigcap_i w_i^{-1}((0, \infty))
  \end{align*}
  as required.
\end{proof}

\begin{lemma}
  \label{lem:int_homothety_cvx}
  \lean{Convex.subset_interior_image_homothety_of_one_lt}\leanok
  Given a point $c$ of $E$ and a real number $t$, let:
  \[
    h^c_t \co E \to E
  \]
  be the homothety which dilates about $c$ by a scale of $t$.

  Suppose $c$ belongs to the interior of a convex subset $C$ of $E$
  and $t > 1$, then
  \[
    C \subseteq \Int(h^c_t(C))
  \]
\end{lemma}

\begin{proof}
  \leanok
  Since $h^c_t$ is a homeomorphism with inverse $h^c_{t^{-1}}$, taking $s = t^{-1}$,
  the required result is equivalent to showing:
  \[
    h^c_s(C) \subseteq \Int(C)
  \]
  where $s \in (0, 1)$.

  Let $x$ be a point of $C$, we must show there exists an open neighborhood $U$
  of $h^c_s(x)$, contained in $C$. In fact we claim:
  \[
    U = h^x_{1-s}(\Int(C))
  \]
  is such a set. Indeed $U$ is open since $h^x_{1-s}$ is a homeomorphism and $U$
  contains $h^c_s(x)$ since:
  \[
    h^c_s(x) = h^x_{1-s}(c) \in h^x_{1-s}(\Int(C))
  \]
  since $c$ belongs to $\Int(C)$. Finally:
  \[
    h^x_{1-s}(\Int(C)) \subseteq h^x_{1-s}(C) \subseteq C
  \]
  where the second inclusion follows since $C$ is convex and contains $x$.
\end{proof}

We are now ready to come back to \Cref{prop:surrounded_by_open}.

\begin{proof}[Proof of \Cref{prop:surrounded_by_open}]\proves{prop:surrounded_by_open}
  \leanok
  \uses{lem:caratheodory,lem:interior_chab,lem:int_homothety_cvx}
  It follows from \Cref{lem:interior_chab} that we need only show that $E$
  has an affine basis $b$ of points belonging to $P$ such that $x$ lies in
  the interior of the convex hull of $b$.

  Carathéodory's lemma~\ref{lem:caratheodory} provides affinely independent
  points $p_0, \dots, p_k$ in $P$ such that $x$ belongs to their convex
  hull. Since $P$ is open, we may extend $p_i$ to an affine basis
  \[
    \hat b = \{p_0, \ldots, p_d\},
  \]
  where all points still belong to $P$.
  Note that $x$ belongs to the convex hull of $\hat b$.

  Now let $c$ be a point in the interior of the convex hull of $\hat b$
  (e.g., the centroid) and for each $\epsilon > 0$, consider the homothety
  \[
    h_{1+\epsilon} \co E \to E,
  \]
  which dilates about $c$ by a scale of $1 + \epsilon$.

  Since $\hat b$ is finite and contained in $P$, and $P$ is open, there exists
  $\epsilon > 0$ such that
  \[
    h_{1+\epsilon} (\hat b) \subseteq P.
  \]
  We claim the required basis is:
  \[
    b = h_{1+\epsilon} (\hat b)
  \]
  for any such $\epsilon$. Indeed, applying
  \Cref{lem:int_homothety_cvx} to $\Conv(\hat b)$ we see:
  \begin{align*}
    x \in \Conv(\hat b) &\subseteq \Int (h_{1+\epsilon} (\Conv(\hat b)))\\
                  &= \Int (\Conv(h_{1+\epsilon} (\hat b)))
  \end{align*}
  as required.
\end{proof}


\section{Constructing loops}

\subsection{Surrounding families}
\label{sub:surrounding_families}

It will be convenient to introduce some more vocabulary.

\begin{definition}
  \label{def:surrounds}
  \uses{def:loop, def:surrounds_points}
  \lean{Loop.Surrounds}
  \leanok
  We say a loop $γ$ surrounds a vector $v$ if $v$ is surrounded
  by a collection of points belonging to the image of $γ$.
  Also, we fix a base point $0$ in $𝕊^1$ and say a loop is based at some
  point $b$ if $0$ is sent to $b$.
\end{definition}

The first main task in proving \Cref{prop:∃_loops} is to construct
suitable families of loops $γ_x$ surrounding $g(x)$, by assembling local
families of loops.
Those will then be reparametrized to get the correct average in the next
section.
In this section, we will work only with \emph{continuous} loops.
This will make constructions easier and we will smooth those loops
in the end, taking advantage of the fact that $Ω$ and the surrounding
condition are open.

Thanks to Carathéodory's lemma, constructing \emph{one} such loop
with values in some open $O$ is easy as soon as $v$ belongs to the
convex hull of $O$.

\begin{lemma}
  \label{lem:loop_of_hull}
  \uses{def:surrounds}
  \lean{surrounding_loop_of_convexHull}
  \leanok
  If a vector $v$ is in the convex hull of a connected open subset $O$
  then, for every base point $b ∈ O$, there is a continuous
  family of loops
  $γ \co [0, 1] × 𝕊^1 → E, (t, s) ↦ γ^t(s)$ such that, for all $t$ and
  $s$:
  \begin{itemize}
    \item
      $γ^t$ is based at $b$
    \item
      $γ^0(s) = b$
    \item
      $γ^t(s) ∈ O$
    \item
      $γ^1$ surrounds $v$
  \end{itemize}
\end{lemma}

\begin{proof}
  \uses{prop:surrounded_by_open}
  \leanok
  Since $O$ is open, \Cref{prop:surrounded_by_open} gives points $p_i$ in $O$
  surrounding $x$.
  Since $O$ is open and connected, it is path connected.
  Let $λ \co [0, 1] → Ω_x$ be a continuous path starting at $b$ and
  going through the points $p_i$.
  We can concatenate $λ$ and its opposite to get $γ^1$.
  This is a round-trip loop: it back-tracks when it reaches $λ(1)$
  at $s = 1/2$.
  We then define $γ^t$ as the round-trip that stops at $s = t/2$, stays
  still until $s = 1-t/2$ and then backtracks.
\end{proof}


\begin{definition}
  \label{def:family_surrounds}
  \uses{def:surrounds}
  \lean{SurroundingFamily}
  \leanok
  A continuous family of loops $γ \co E × [0, 1] × 𝕊^1 → F, (x, t, s) ↦ γ^t_x(s)$
  surrounds a map $g \co E → F$ with base $β \co E → F$
  on $U ⊂ E$ in $Ω ⊂ E × F$ if, for every $x$ in $U$, every $t ∈ [0, 1]$ and every
  $s ∈ 𝕊^1$,
  \begin{itemize}
    \item
      $γ^t_x$ is based at $β(x)$
    \item
      $γ^0_x(s) = β(x)$
    \item
      $γ^1_x$ surrounds $g(x)$
    \item
      $(x,γ^t_x(s)) ∈ Ω$.
  \end{itemize}
  The space of such families will be denoted by
  $\Loop(g, β, U, Ω)$.
\end{definition}

Families of surrounding loops are easy to construct locally.

\begin{lemma}
  \label{lem:local_loops}
  \uses{def:family_surrounds}
  \lean{local_loops}
  \leanok
  Assume $Ω$ is open over some neighborhood of $x_0$.
  If $g(x_0)$ is in the convex hull of the connected component of $Ω_{x_0}$ containing $β(x_0)$,
  then there is a continuous family of loops defined near $x_0$, based at $β$,
  taking value in $Ω$ and surrounding $g$.
\end{lemma}

\begin{proof}
  \leanok
  \uses{lem:loop_of_hull, lem:smooth_barycentric_coord}
  In this proof we don't mention the $t$ parameter since it plays no
  role, but it is still there.
  \Cref{lem:loop_of_hull} gives a loop $γ$ based at $β(x_0)$,
  taking values in $Ω_{x_0}$ and surrounding $g(x_0)$.
  We set $γ_x(s) = β(x) + (γ(s) - β(x_0))$.
  Each $γ_x$ takes values in $Ω_x$ because $Ω$ is open over some
  neighborhood of $x_0$.
  \Cref{lem:smooth_barycentric_coord} guarantees that this loop surrounds $g(x)$
  for $x$ close enough to $x_0$.
\end{proof}

The difficulty in constructing global families of surrounding loops is that
there are plenty of surrounding loops and we need to choose them
consistently.
The key feature of the above definition is that the $t$ parameter not only
allows us to cut out the corrugation
process in the next chapter, but also brings a ``satisfied or refund'' guarantee,
as explained in the next lemma.

\begin{lemma}
  \label{lem:satisfied_or_refund}
  \uses{def:family_surrounds}
  \lean{satisfied_or_refund}
  \leanok
  For every set $U ⊂ E$,  $\Loop(g, β, U, Ω)$ is ``path connected'':
  for every $γ_0$ and $γ_1$ in $\Loop(g, β, U, Ω)$,
  there is a continuous map
  $δ \co [0, 1] × E × [0, 1] × 𝕊^1 → F, (τ, x, t, s) ↦ δ^t_{τ, x}(s)$
  which interpolates between $γ_0$ and $γ_1$ in
  $\Loop(g, β, U, Ω)$.
\end{lemma}

The construction below morally proves that each $\Loop(g, β, U, Ω)$ is contractible, but
we will not even specify a topology on those spaces.
The definition of ``path connected'' in quotation marks is the above specific statement,
and only this statement will be used.

\begin{proof}
  \leanok
  Let $ρ$ be the piecewise affine map from $ℝ$ to $ℝ$ such that
  $ρ(τ) = 1$ if $τ ≤ 1/2$, $ρ$ is affine on $[1/2, 1]$,
  $ρ(τ) = 0$ if $τ ≥ 1$.
  We set
  \[
    δ_{τ, x}^t(s) =
    \begin{cases}
      γ_{0,x}^{ρ(τ)t}\left(\frac1{1 - τ} s\right) & \text{if $s ≤ 1 - τ$ and $τ < 1$}\\
      γ_{1,x}^{ρ(1-τ)t}\left(\frac1τ \big(s - (1- τ)\big)\right) &
      \text{if $s ≥ 1 - τ$ and $τ > 0$}\\
    \end{cases}
  \]
  It is clear that if $s = 1 - τ$ then both branches agree and are equal to $β(x)$.
  Therefore it is easy to see that $δ$ is continuous at $(τ, x, t, s)$
  except when $(τ,s)=(1,0)$ or $(τ,s)=(0,1)$.

  To show the continuity for $(τ,s)=(1,0)$, let $K$ be a compact neighborhood of $x$ in $E$.
  Then $γ_0$ is uniformly continuous on the compact set $K × [0, 1] × 𝕊^1$, which means that
  $γ_{0,x'}^t$ tends uniformly to the constant function $s ↦ β(x)$ as $(x', t)$ tends to
  $(x, 0)$.
  This means that $γ_{0,x'}^{ρ(τ)t'}$ tends uniformly to the constant function $s ↦ β(x)$
  as $(τ, x', t')$ tends to $(1, x, t)$. This means that $δ$ is continuous at $(τ,s)=(1,0)$
  (it is clear that the other branch also tends to $β(x)$). The continuity at $(τ,s)=(0,1)$ is
  entirely analogous.

  The beautiful observation motivating the above formula is why each
  $δ_{τ, x}^1$ surrounds $g(x)$.
  The key is that the image of $δ_{τ, x}^1$ contains the image of
  $γ_{0,x}^1$ when $τ ≤ 1/2$, and contains  the image of
  $γ_{1,x}^1$ when $τ ≥ 1/2$.
  Hence $δ_{τ, x}^1$ always surrounds $g(x)$.
\end{proof}

\begin{corollary}
  \label{cor:extend_loops}
  % \uses{def:family_surrounds} % causes transitive arrow in depgraph
  \lean{extend_loops}
  \leanok
  Let $U_0$ and $U_1$ be open sets in $E$.
  Let $K_0 ⊂ U_0$ and $K_1 ⊂ U_1$ be compact subsets.
  For any $γ_0 ∈ \Loop(U_0, g, β, Ω)$ and $γ_1 ∈ \Loop(U_1, g, β, Ω)$,
  there exists a neighborhood $U$ of $K_0 ∪ K_1$ and
  there exists $γ ∈ \Loop(U, g, β, Ω)$
  which coincides with $γ_0$ near $K_0\cup U_1^c$.
\end{corollary}

\begin{proof}
  \leanok
  \uses{lem:satisfied_or_refund}
  Let $C_0 = K_0\cup U_1^c$ and $C_1 := K_1 ∖ U_0$. Since $C_0$ and $C_1$ are disjoint closed sets,
  there is some continuous cut-off $ρ \co E → [0, 1]$
  which vanishes on a neighborhood of $C_0$ and equals one on a neighborhood of $C_1$.

  \Cref{lem:satisfied_or_refund} gives a homotopy of loops
  $γ_τ$ from $γ_0$ to $γ_1$ on $U_0 ∩ U_1$. Moreover, note that $γ_τ$ is defined on all of $E$.
  On $U_0' ∪ (U_0 ∩ U_1) ∪ U_1'$, which is a
  neighborhood of $K_0 ∪ K_1$, we set
  \[
      γ_x = γ_{ρ(x), x}
  \]
  which has the required properties.
\end{proof}

\begin{lemma}
  \label{lem:∃_surrounding_loops}
  % \uses{def:family_surrounds} % causes transitive arrow in depgraph
  \lean{exists_surrounding_loops}
  \leanok
  In the setup of \Cref{prop:∃_loops}, assume we have a
  continuous family $γ$ of loops defined near $K$ which is based at $β$,
  surrounds $g$ and such that each $γ_x^t$ takes values in $Ω_x$.
  Then there such a family which is defined on all of $E$ and agrees
  with $γ$ near $K$.
\end{lemma}

\begin{proof}
  \uses{lem:local_loops, cor:extend_loops, lem:relative_inductive_construction_of_loc}
  \leanok
  \Cref{lem:local_loops} proves the existence of local families of surrounding
  loops and \Cref{cor:extend_loops} allows to patch such families hence
  \Cref{lem:relative_inductive_construction_of_loc} proves global existence.
\end{proof}

\subsection{The reparametrization lemma}
\label{sub:the_reparametrization_lemma}

The second ingredient needed to prove \Cref{prop:∃_loops} is a
parametric reparametrization lemma. Gromov's original proof of this lemma makes
explicit use of a partition of unity. Motivated in particular by formalization
purposes, we will first state more abstract versions whose statements do not involve
any partition of unity but directly state a local-to-global property.

\begin{lemma}
  \label{lem:exists_cont_diff_of_convex}
  \lean{exists_contDiff_of_convex}
  \leanok
  Let $E$ and $F$ be real normed vector spaces. Assume that $E$ is finite
  dimensional. Let $P$ be a predicate on $E \times F$ such that for all $x$ in
  $E$, $\{y ~|~ P (x, y) \}$ is convex. Let $n$ be a natural number or $+\infty$.
  Assume that every $x$ has a neighbourhood $U$ on which there exists a $\C^n$
  function $f$ such that $\forall x ∈ U, P(x, f(x))$. Then there is a global
  $\C^n$ function $f$ such that $\forall x, P(x, f(x))$.
\end{lemma}

\begin{proof}
  \leanok
  The assumption give us an open cover $(U_i)_{i ∈ I}$ of $E$ and functions
  $f_i \co E → F$ that are smooth on $U_i$ and such that $P(x, f_i(x))$ for all $x$ in $U_i$.
  Let $ρ$ be a smooth partition of unity associated to this cover. The function
  $f = ∑ ρ_i f_i$ is smooth on $E$ and the convexity assumption on $P$ ensures
  it satisfies $\forall x, P(x, f(x))$. Indeed each value $f(x)$ is a convex combination
  of finitely many values $f_i(x)$ where $i$ satisfies that $x$ is in $U_i$.
\end{proof}

We will also need a version where $F$ is a space of smooth functions. Since
there is no relevant norm to put on such a space, we cannot deduce this version
from the above one.

\begin{lemma}
  \label{lem:exists_cont_diff_of_convex₂}
  \lean{exists_contDiff_of_convex₂}
  \leanok
  Let $E₁$, $E₂$ and $F$ be real vector spaces. Assume $E₁$ and $E₂$ are finite dimensional.
  Let $n$ be a natural number or $+\infty$. Let $P$ be a property
  of pairs $(x, f)$ with $x ∈ E₁$ and $f : E₂ → F$. Assume that, for every $x$, the space
  of functions $f$ such that $P(x, f)$ holds is convex. Assume that for every $x₀$
  in $E₁$ there is a neighborhood $U$ of $x₀$ and a function $φ : E₁ × E₂ → F$ which is
  $\C^n$ on $U × E₂$ and such that $P(x, φ(x, \cdot))$ holds for every $x$ in $U$.
  There there is a global $\C^n$ function $φ \co E₁ × E₂ → F$ such that
  $P(x, φ(x, \cdot))$ holds for every $x$.
\end{lemma}

\begin{proof}
  \leanok
  This is completely analogous to the previous proof.
\end{proof}

\begin{lemma}
\label{lem:reparametrization}
\uses{def:surrounds}
\lean{SmoothSurroundingFamily.reparametrize_average}
\leanok
Let $γ \co E × 𝕊^1 → F$ be a smooth family of loops surrounding
a map $g$.
There is a smooth family $φ \co E × 𝕊^1 → 𝕊^1$ such
that each $γ_x ∘ φ_x$ has average $g(x)$ and $φ_x(0) = 0$.
\end{lemma}

\begin{proof}
  \uses{lem:smooth_barycentric_coord, lem:exists_cont_diff_of_convex₂}
  \leanok
  Gromov's main idea in order to prove this result is to translate the problem
  of constructing a family of circle maps $φ$ into the problem of constructing
  a family of smooth density functions $f$ on the circle. We introduce some
  vocabulary in order to describe this reduction.
  Let $f \co E × ℝ → ℝ$ be a smooth positive function that is $1$-periodic in
  its second argument.
  We say that $f$ is a centering density for $(γ, g)$ at $x$ if
  $f_x \co ℝ → ℝ$ has average value one when seen as a function on $𝕊¹$ and
  the average value of $f_x γ_x$ is $g(x)$. We claim that, in order to prove
  the lemma, it is sufficient to build such an $f$ which is centering at every
  $x$. Indeed, assume we have such an $f$. We then get a smooth family
  of $ℤ$-equivariant functions $ψ \co E × ℝ → ℝ$ defined by
  $ψ_x(t) = \int_0^tf_x(s)ds$. Because $ψ$ is smooth and each $ψ_x$ is strictly
  monotone and $ℤ$-equivariant, one can check there is a smooth map
  $φ : E × ℝ → ℝ$ which is $ℤ$-equivariant and such that $φ_x ∘ ψ_x = \Id$ for
  each $x$. Seen as a family of functions from $𝕊¹$ to $𝕊¹$, those functions
  are suitable since, for every $x$, the change of variable formula gives:
  \[
    \int_{𝕊¹} γ_x ∘ φ_x(s)ds = \int_{𝕊¹} ψ_x'(s) γ_x ∘ φ_x(ψ_x(s))ds
     = \int_{𝕊¹} f_x(s) γ_x(s)ds = g(x).
  \]

  We now prove the existence of a function which is a centering density at every point of $x$.
  For any given $x$, this constraint is clearly convex.
  Hence \Cref{lem:exists_cont_diff_of_convex₂} ensures it is enough to prove existence
  of functions that are centering densities in a neighborhood of any given point $x$.
  So we fix some $x$ in $E$.

  Since $γ_x$ strictly surrounds $g(x)$, there are points
  $s_1, …, s_{n+1}$ in $𝕊^1$ such that $g(x)$ is surrounded
  by the corresponding points $γ_x(s_j)$.


  Let $f_1, …, f_{n+1}$ be smooth positive periodic maps from $ℝ$ to $ℝ$
  which average value $1$ on a period and such that the corresponding measures
  on $𝕊¹$ are very
  close to the Dirac measures on $s_j$, ie. for any function $h$, the average value of
  $f_jh$ is almost $h(s_j)$. We set $p_j = \int f_jγ_x\, ds$, which
  is almost $γ_x(s_j)$ so that $g(x) = \sum w_j p_j$ for some weights
  $w_j$ in the open interval $(0, 1)$ according to
  \Cref{lem:smooth_barycentric_coord}.

  If $x'$ is in a sufficiently small neighborhood of $x$,
  \Cref{lem:smooth_barycentric_coord} gives smooth weight functions $w_j$
  such that $g(x') = \sum w_j(x')p_j(x')$. Hence we can set
  $f_{x'}(s) = \sum w_j(x')f_j(s)$.
\end{proof}

\subsection{Proof of the loop construction proposition}
\label{sub:proof_of_the_loop_construction_proposition}

We finally assemble the ingredients from the previous two sections.

\begin{proof}[Proof of \Cref{prop:∃_loops}]
  \proves{prop:∃_loops}
  \uses{lem:∃_surrounding_loops, lem:reparametrization, lem:exists_cont_diff_of_convex}
  \leanok
  Let $γ^*$ be a family of loops surrounding the origin in $B_F(0,1)$
  the open unit ball in $F$, constructed using \Cref{lem:local_loops}.
  For $x$ in some neighborhood $U^*$ of $K$ where $g = β$, we set
  $γ_x = g(x) + εγ^*$ where $ε > 0$ is sufficiently small to ensure that
  $B_{E\times F}((x,β(x)),2ε)\subseteq \Omega$ (recall $Ω$ is
  open and $K$ is compact).
  \Cref{lem:∃_surrounding_loops} extends this family to a continuous
  family of surrounding loops $γ_x$ for all $x$ (this is not yet our
  final $γ$).

  We then need to approximate this continuous family by a smooth one.
  Some care is needed to ensure that it stays based at $β$.
  We can first reparametrize $γ$ on $[0,1] \times 𝕊^1$
  to ensure that $γ$ is constant in a neighborhood of
  $C = \{(t, s) \in [0,1] \times 𝕊^1 \mid t = 0 \text{ or } s = 0\}$.
  Using \Cref{lem:exists_cont_diff_of_convex}, we can find a smooth function
  that has distance at most $ε$ from $γ$ and coincides with $γ$ on $C$
  (using the fact that $γ$ is already smooth near $C$).
  Since all loops that are sufficiently close to $γ$ still surround $g$,
  we can also ensure that the new smoothened $γ$ is still surrounding.

  Then \Cref{lem:reparametrization} gives a family of circle
  diffeomorphisms $h_x$ such that $γ^1_x ∘ h_x$ has average $g(x)$.

  Finally we choose a cut-off function function $χ$ which vanishes near
  $E ∖ U^*$ and equals one near $K$. As our final family of loops,
  we choose $χ(x)g(x) + (1-χ(x))(γ_x ∘ h_x)$. This
  operation does not change the average values of these loops, because
  it rescales them around their average value, but makes them constant
  near $K$. Also, those loops stay in $Ω$, thanks to our choice of $ε$.
\end{proof}

% vim: set expandtab sw=2 tabstop=2:

\chapter{Local theory of convex integration}
\label{chap:local}
\section{Key construction}
\label{sec:convex_integration_intro}

The goal of this chapter is to explain the local aspects of
(Theillière's implementation of) convex integration, the next chapter will
cover global aspects.

The elementary step of convex integration modifies the derivative of a
map in one direction. The precise meaning of ``one direction'' relies on the
following definition.

\begin{definition}
  \label{def:dual_pair}
  \leanok
  \lean{DualPair}
  A dual pair on a vector space $E$ is a pair $(π, v)$ where $π$ is a linear
  form on $E$ and $v$ a vector in $E$ such that $π(v) = 1$.
\end{definition}

Let $E$ and $F$ be finite dimensional real normed vector spaces.
Let $f \co E → F$ be a smooth map, and let $(π, v)$ be a dual pair on $E$.
We want to modify $Df$ in the direction of $v$ while almost preserving it
on $\ker π$.
Say we wish $Df(x)v$ could live in some open subset $Ω_x ⊂ F$. Assume there is
a smooth family of loops $γ \co E × 𝕊^1 → F$ such that each
$γ_x$ takes values in $Ω_x$, and its average value $\overline{γ}_x = \int_{𝕊^1}
γ_x$ is $Df(x)v$ for all $x$. Obviously such loops can exist
only if $Df(x)v$ is in the convex hull of $Ω_x$, and we saw in the previous
chapter that this is almost sufficient (and we'll see this is sufficiently
almost sufficient for our purposes). Then we can modify $f$ to fulfil our wish
using the following construction.

\begin{definition}[Theillière 2018]
  \label{def:corrugation}
  \uses{def:loop}
  \leanok
  \lean{corrugation}
  The map obtained by corrugation of $f$ in direction $(π, v)$ using $γ$
  with oscillation number $N$ is
  \[
    x ↦ f(x) + \frac1N  ∫_0^{Nπ(x)} \left[γ_x(s) - \overline{γ}_x\right]ds.
  \]
\end{definition}

In the above definition, we mostly think of $N$ as a large natural
number.
But we don't actually require it, any positive real number will do.

The next proposition implies that, provided $N$ is large enough, we have
achieved $Df'(x)v ∈ Ω_x$, almost without modifying derivatives in the
directions of $\ker π$, and almost without moving $f(x)$.

\begin{proposition}[Theillière 2018]
  \label{prop:theilliere}
  \lean{fderiv_corrugated_map, corrugation.c0_small_on, remainder_c0_small_on, corrugation.support}\leanok
  \uses{def:corrugation}
  Let $f$ be a $\mathcal{C}^1$ function from $E$ to $F$. Let $(π, v)$ be a dual
  pair on $E$. Let $γ \co E × 𝕊^1 → F$ be a $\mathcal{C}^1$ family of loops
  such that $\overline{γ_x} = Df(x)v$ for all $x$.

  For any compact set $K ⊂ E$ and any positive $ε$, the function $f'$ obtained
  by corrugation of $f$ in direction $(π, v)$ using $γ$ with large enough
  oscillation number $N$ satisfies:
  \begin{enumerate}
    \item\label{CP:C0}
      $∀ x ∈ K, ‖f'(x) - f(x)‖ ≤ ε$
    \item\label{CP:kerpi}
      $∀ x ∈ K, ‖(Df'(x) - Df(x))_{|\ker π}‖ ≤ ε$.
    \item\label{CP:v}
      $∀ x ∈ K, ‖Df'(x)v - γ(x, Nπ(x))‖ ≤ ε$
  \end{enumerate}
  In addition, all the differences estimated above vanish if $x$ is outside the
  support of $γ$.
\end{proposition}

\begin{proof}
  \leanok
  We set $Γ_x(t) = ∫_0^t \left(γ_x(s) - \overline{γ}_x\right)ds$, so that
  $f'(x) = f(x) + Γ_x(Nπ(x))/N$.  Because each $Γ_x$ is $1$-periodic, and
  everything has compact support in $E$, all derivatives of $Γ$ are uniformly
  bounded.  Item~\ref{CP:C0} in the statement is then obvious.
  Item~\ref{CP:kerpi} also follows since
  $∂_if'(x) = ∂_if(x) + ∂_iΓ(x, Nπ(x))/N$. In order to prove
  Item~\ref{CP:v}, we compute:
  \begin{align*}
      Df'(x)v &= Df(x)v + \frac1N ∂_jΓ(x, Nπ(x)) + \frac NN ∂_tΓ(x, Nπ(x))\\
                       &= Df(x)v + \bigO{\frac1N} + γ(x, Nπ(x)) - Df(x)v\\
                       &= γ(x, Nπ(x)) + \bigO{\frac1N}.
  \end{align*}
  Outside the support of $γ$, $Γ_x$ and its derivative with respect to $x$
  vanish identically (for the derivative computation, it is important that the
  support of $γ$ is the \emph{closure} of the set of $x$ where $γ_x$ is not
  constant).
\end{proof}

\section{The main inductive step}
\label{sec:inductive_step}

\begin{definition}
  \label{def:hol_partial}
  \leanok
  \lean{JetSec.IsPartHolonomicAt}
  Let $E'$ be a linear subspace of $E$.
  A map $\F = (f, φ) : E → F × \Hom(E, F)$ is $E'$--holonomic if,
  for every $v$ in $E'$ and every $x$, $Df(x)v = φ(x)v$.
\end{definition}

\begin{definition}
  \label{def:rel_loc}
  \leanok
  \lean{RelLoc}
  A first order differential relation for maps from $E$ to $F$ is a
  subset $\Rel$ of $E × F × \Hom(E, F)$.
\end{definition}

Until the end of this section, $\Rel$ will always denote a first order
differential relation for maps from $E$ to $F$.

\begin{definition}
  \label{def:formal_sol_loc}
  \leanok
  \lean{RelLoc.FormalSol}
  \uses{def:rel_loc}
  A formal solution of a differential relation $\Rel$ is a map
  $\F = (f, φ) \co E → F × \Hom(E, F)$ such that, for every $x$,
  $(x, f(x), φ(x))$ is in $\Rel$.
\end{definition}

The first component of a map $\F \co E → F × \Hom(E, F)$ will sometimes
be denoted by $\bs \F \co E → F$ and called the base map of $\F$.

\begin{definition}
  \label{def:htpy_jet_sec_loc}
  \leanok
  \lean{HtpyJetSec}
  A $1$-jet section from $E$ to $F$ is a function from $E$ to $F × \Hom(E, F)$.
  A homotopy of $1$-jet sections is a smooth map
  $\F : ℝ × E → F × \Hom(E, F)$.
\end{definition}

Typically, $x ↦ \F(t, x)$ will be denoted by $\F_t$. It could seem more natural to take
$[0, 1] × E$ as the source of a homotopy but this would be less convenient for
formalization and wouldn't change anything since any map from $ℝ × E$ can be restricted to
$[0, 1] × E$ and every map from $[0, 1] × E$ could be extended.

\begin{definition}
  \label{def:rel_slice}
  \leanok
  \lean{RelMfld.slice}
  \uses{def:rel_loc}
  For every $σ = (x, y, φ)$, the slice of $\Rel$ at $σ$ with respect to $(π, v)$
  is:
  \[
    \Rel(σ, π, v) = \{w ∈ F \;|\; (x, y, φ + (w - φ(v)) ⊗ π) ∈ \Rel\}.
  \]
\end{definition}

\begin{lemma}
  \label{lem:update_lin_map}
  \leanok
  \lean{DualPair.update_ker_pi, DualPair.update_v, mem_slice}
  The linear map $φ + (w - φ(v)) ⊗ π)$ coincides with $φ$ on $\ker π$
  and sends $v$ to $w$.
  If $\sigma$ belongs to $\Rel$ then $φ(v)$ belongs to
  $\{w ∈ F, (x, y, φ + (w - φ(v)) ⊗ π) ∈ \Rel\}$.
\end{lemma}

\begin{proof}\leanok
  These are direct checks.
\end{proof}


We'll use the notation $\Conn_w A$ to denote the connected component of $A$
that contains $w$, or the empty set if $w$ doesn't belong to $A$.

\begin{definition}
  \label{def:short_formal_sol}
  \leanok
  \lean{RelLoc.FormalSol.IsShortAt}
  \uses{def:formal_sol_loc, def:dual_pair}
  A formal solution $\F$ of $\Rel$ is $(π, v)$--short if,
  for every $x$, $Df(x)v$ belongs to the interior of the convex hull of
  $\Conn_{φ(v)}\Rel((x, f(x), φ(x)), π, v)$.
\end{definition}

\begin{lemma}
  \label{lem:integration_step}
  \leanok
  \lean{StepLandscape.improveStep_rel_t_eq_0, StepLandscape.improveStep_rel_C,%
  StepLandscape.improveStep_rel_compl_K₁, StepLandscape.improveStep_c0_close,%
  StepLandscape.improveStep_part_hol}
  \uses{def:short_formal_sol, def:hol_partial, def:htpy_jet_sec_loc}
  Let $\F$ be a formal solution of $\Rel$.
  Let $K_1 ⊂ E$ be a compact subset, and let $K_0$ be a compact subset of
  the interior of $K_1$. Let $C$ be a closed subset of $E$.
  Let $E'$ be a linear subspace of $E$ contained in $\ker π$.
  Let $ε$ be a positive real number.

  Assume $\Rel$ is open.
  Assume that $\F$ is $E'$--holonomic near $K_0$, $(π, v)$--short,
  and holonomic near $C$.
  Then there is a homotopy $\F_t$ such that:
  \begin{enumerate}
    \item
      $\F_0 = \F$~;
    \item
      $\F_t$ is a formal solution of $\Rel$ for all $t$~;
    \item
      $\F_t(x) = \F(x)$ for all $t$ when $x$ is near $C$ or outside
      $K_1$~;
    \item
      $d(\bs\F_t(x), \bs \F(x)) ≤ ε$ for all $t$ and all $x$~;
    \item
      $\F_1$ is $E' ⊕ ℝv$--holonomic near $K_0$.
  \end{enumerate}
\end{lemma}

\begin{proof}
  \leanok
  \proves{lem:integration_step}
  \uses{prop:∃_loops, prop:theilliere, lem:update_lin_map}
  We denote the components of $\F$ by $f$ and $φ$.
  Since $\F$ is short, \cref{prop:∃_loops} applied to
  $g \co x ↦ Df(x)v$, $β \co x ↦ φ(x)v$,
  $Ω_x = \Rel(\F(x), π, v)$, and $K = C ∩ K_1$ gives us a smooth family of loops
  $γ \co E × [0, 1] × 𝕊^1 → F$ such that, for all $x$:
  \begin{itemize}
    \item $∀ t\, s,\; γ^t_x(s) ∈ \Rel(\F(x), π, v)$
    \item $∀ s,\; γ^0_x(s) = φ(x)v$
    \item $\bar γ^1_x = Df(x)v$
    \item if $x$ is near $C$, $∀t\, s,\; γ^t_x(s) = φ(x)v$
  \end{itemize}
  Let $ρ: E → ℝ$ be a smooth cut-off function which equals one on
  a neighborhood of $K_0$ and whose support is contained in $K_1$.

  Let $N$ be a positive real number.
  Let $\bar f$ be the corrugated map constructed from $f$, $γ^1$ and $N$.
  \Cref{prop:theilliere} ensures that, for all $x$,
  \[
    D\bar f(x) = Df(x) + \left[γ^1_x(Nπ(x)) - Df(x)v\right] ⊗ π +
             R_x
  \]
  for some remainder term $R$ which is $ε$-small and vanishes whenever $γ_x$ is
  constant, hence vanishes near $C$.

  We set $\F_t(x) = \big(f_t(x), φ_t(x)\big)$ where:
  \[
    f_t(x) = f(x) + \frac{tρ(x)}N \int_0^{Nπ(x)}
                    \left[γ^t_x(s) - Df(x)v\right]\, ds
  \]
  and
  \[
    φ_t(x) = φ(x) + \left[γ^{tρ(x)}_x(Nπ(x)) - φ(x)v\right] ⊗ π +
            \frac{tρ(x)}N B_x.
  \]

  We now prove that $\F_t$ has the announced properties, starting with
  he obvious ones.
  The fact that $\F_0 = \F$ is obvious since $γ^0_x(s) = φ(x)v$ for all
  $s$.

  When $x$ is near $C$, $Df(x) = φ(x)$ since $\F$ is holonomic near
  $C$.
  In addition, $γ^t_x(s) = φ(x)v$ for all $s$ and $t$, hence $B_x$ vanishes.
  Hence $\F_t(x) = \F(x)$ for all $t$ when $x$ is near $C$.

  Outside of $K_1$, $ρ$ vanishes. Hence $f_t(x) = f(x)$ for all $t$,
  and $γ^{tρ(x)}_x(s) = φ(x)v$ for all $s$ and $t$, and $φ_t(x) = φ(x)$.

  The distance between $f(x)$ and $f_t(x)$ is zero outside of $K_1$
  and $ε$-small everywhere.

  We now turn to the interesting parts.
  The first one is that each $\F_t$ is a formal solution of $\Rel$.
  We already now that $\F_t$ coincides with $\F$, which is a formal
  solution, outside of the compact set $K_1$.
  We set
  \[
    \F'_t(x) = \left(f(x),
                     φ(x) + \left[γ^{tρ(x)}_x(Nπ(x)) - φ(x)v\right] ⊗ π\right).
  \]
  Since $\Rel$ is open, and $K_1 × [0, 1]$ is compact and $\F_t$ is
  within $\bigO{1/N}$ of $\F'_t$, it suffices to prove that $\F'_t$ is a formal
  solution for all $t$.
  This is guaranteed by the definition of the slice $\Rel(\F(x), π, v)$
  to which $γ^{tρ(x)}_x(Nπ(x))$ belongs.

  Finally, let's prove that $\F_1$ is $E' ⊕ ℝv$--holonomic near $K_0$.
  Since $ρ = 1$ near $K_0$, we have, for $x$ near $K_0$,
  \[
    Df_1(x) = Df(x) + \left[γ^1_x(Nπ(x)) - Df(x)v\right] ⊗ π +
             \frac1N B_x,
  \]
  and
  \[
    φ_1(x) = φ(x) + \left[γ^1_x(Nπ(x)) - φ(x)v\right] ⊗ π +
            \frac{1}N B_x.
  \]
  Let $p$ be the projection of $E$ onto $\ker π$ along $v$,
  so that $\Id_E = p + v ⊗ π$.
  We can rewrite the above formulas as
  \[
    Df_1(x) = Df(x) ∘ p + γ^1_x(Nπ(x)) ⊗ π + \frac1N B_x,
  \]
  and
  \[
    φ_1(x) = φ(x) ∘ p + γ^1_x(Nπ(x)) ⊗ π + \frac{1}N B_x.
  \]
  So we see the difference is $Df(x) ∘ p - φ(x) ∘ p$ which vanishes on
  $E'$ since $\F$ is $E'$--holonomic near $K_0$, and vanishes on $v$ since
  $p(v) = 0$.
\end{proof}

\section{Ample differential relations}
\label{sec:ample_differential_relations}

\begin{definition}
  \label{def:ample_subset}
  \leanok
  \lean{AmpleSet}
  A subset $Ω$ of a real vector space $E$ is ample if the convex
  hull of each connected component of $Ω$ is the whole $E$.
\end{definition}

\begin{lemma}
  \label{lem:ample_codim_two}
  \uses{def:ample_subset}
  \lean{AmpleSet.of_one_lt_codim}
  \leanok
  The complement of a linear subspace of codimension at least 2 is
  ample.
\end{lemma}

\begin{proof}
  \leanok
  Let $F$ be subspace of $E$ with codimension at least $2$.
  Let $F'$ be a complement subspace.
  Its dimension is at least $2$ since it is isomorphic to $E/F$
  and $\dim(E/F) = \codim(F) ≥ 2$.
  First note the complement of $F$ is path-connected.
  Indeed let $x$ and $y$ be points outside $F$.
  Decomposing on $F ⊕ F'$, we get $x = u + u'$ and $y = v + v'$
  with $u' ≠ 0$ and $v' ≠ 0$.
  The segments from $x$ to $u'$ and $y$ to $v'$ stay outside $F$,
  so it suffices to connect $u'$ and $v'$ in $F' ∖ \{0\}$.
  If the segment from $u'$ to $v'$ doesn't contains the origin then we
  are done.
  Otherwise $v' = μu'$ for some (negative) $u'$.
  Since $\dim(F') ≥ 2$ and $u' ≠ 0$, there exists $f ∈ F'$
  which is linearly independent from $u'$, hence from $v'$.
  We can then connect both $u'$ and $v'$ to $f$ by a segment
  away from zero.

  We now turn to ampleness.
  The connectedness result reduces to prove that every $e$ in $E$
  is in the convex hull of $E ∖ F$.
  If $e$ is not in $F$ then it is the convex combination of itself with
  coefficient $1$ and we are done.
  Now assume $e$ is in $F$.
  The codimension assumption guarantees the existence
  of a subspace $G$ such that $\dim(G) = 2$ and $G ∩ F = \{0\}$.
  Let $(g_1, g_2)$ be a basis of $G$.
  We set $p_1 = e + g_1$, $p_2 = e + g_2$, $p_3 = e - g_1 - g_2$.
  All these points are in $E ∖ F$ and $e = p_1/3 + p_2/3 + p_3/3$.
\end{proof}

\begin{definition}
  \label{def:ample_relation_loc}
  \leanok
  \lean{RelLoc.IsAmple}
  \uses{def:ample_subset, def:rel_slice}
  A first order differential relation $\Rel$ is ample if all its slices
  are ample.
\end{definition}

\begin{lemma}
  \label{lem:h_principle_open_ample_loc}
  \leanok
  \lean{RelLoc.FormalSol.improve}
  \uses{def:ample_relation_loc, def:holonomic_loc}
  Let $\F$ be a formal solution of $\Rel$.
  Let $K_1 ⊂ E$ be a compact subset, and let $K_0$ be a compact subset of
  the interior of $K_1$.
  Assume $\F$ is holonomic near a closed subset $C$ of $E$.
  Let $ε$ be a positive real number.

  If $\Rel$ is open and ample then there is a homotopy $\F_t$ such that:
  \begin{enumerate}
    \item
      $\F_0 = \F$
    \item
      $\F_t$ is a formal solution of $\Rel$ for all $t$~;
    \item
      $\F_t(x) = \F(x)$ for all $t$ when $x$ is near $C$ or outside $K_1$.
    \item
      $d(\bs\F_t(x), \bs \F(x)) ≤ ε$ for all $t$ and all $x$~;
    \item
      $\F_1$ is holonomic near $K_0$~;
    \item
      $t ↦ F_t$ is constant near $0$ and $1$.
  \end{enumerate}
\end{lemma}

\begin{proof}
  \leanok
  \uses{lem:integration_step}
  This is a straightforward induction using \cref{lem:integration_step}.
  Let $(e_1, \dots, e_n)$ be a basis of $E$,
  and let $(π_1, \dots, π_n)$ be the dual basis.
  Let $E'_i$ be the linear subspace of $E$ spanned by
  $(e_1, \dots, e_i)$, for $1 ≤ i ≤ n$, and let $E'_0$ be the zero
  subspace of $E$.
  Each $(π_i, e_i)$ is a dual pair and the kernel of $π_i$ contains
  $E'_{i - 1}$.

  \Cref{lem:integration_step} allows to build a sequence of homotopies
  of formal solutions,
  each homotopy relating a formal solution which is $E'_i$--holonomic to
  one which is $E'_{i+1}$--holonomic (always near $K_0$).
  The shortness condition is always satisfies because $\Rel$ is
  ample.
  Each homotopy starts where the previous one stopped, stay at
  $C^0$ distance at most $ε/n$, and is relative to $C$ and
  the complement of $K_1$.

  It then suffices to do a smooth concatenation of theses homotopies.
  We first pre-compose with a smooth map from $[0, 1]$ to itself that
  fixes $0$ and $1$ and has vanishing derivative to all orders at $0$
  and $1$.
  Then we precompose by affine isomorphisms from $[0, 1]$
  to $[i/n, (i + 1)/n]$ before joining them.
\end{proof}


% vim: set expandtab sw=2 tabstop=2:

\chapter{Global theory of open and ample relations}
\label{chap:global}

\section{Preliminaries}

\subsection{Localisation data}%
\label{sub:localisation_data}

In order to conveniently globalize the theory of the previous chapter, we'll
need a number of constructions and lemmas. By definition, manifolds are covered
by open sets that are diffeomorphic to open sets of vector spaces. But for us it is
slightly more convenient to work with smooth open embeddings of whole vector spaces.
Here a smooth open embedding from a manifold $X$ to a manifold $Y$ is a smooth map
$φ : X → Y$ which is open and for which there is some smooth $ψ : φ(X) → X$ such that
$ψ ∘ φ = \Id_X$ (hence also and $φ ∘ ψ = \Id_{φ(X)}$).
Remember that a family of sets $V_i$ in a topological space $X$ is locally finite if
every point of $X$ has a neighborhood that intersects only finitely many $V_i$.
Note that in this whole text, every manifold is paracompact by definition. In
particular their topology are metrizable and we will arbitrarily fix a
compatible distance function on every manifold.

\begin{definition}
  \label{def:update}
  \lean{OpenSmoothEmbedding.update}
  \leanok
  Given smooth open embeddings $φ : X → M$ and $ψ : Y → N$, the update of a map
  $f : M → N$, using a map $g : X → Y$
  is the map from $M$ to $N$ sending $m$ to $ψ ∘ g ∘ φ⁻¹(m)$ if
  $m ∈ φ(X)$ and $f(m)$ otherwise.
\end{definition}

\begin{lemma}
  \label{lem:smooth_updating}
  \uses{def:update}
  \lean{OpenSmoothEmbedding.contMDiff_update}
  \leanok
  Let $φ : P × X → M$ and $ψ : P × Y → N$ be families of smooth open embeddings.
  Let $K$ be a set in $X$ whose image in $M$ is closed. Let $f : P × M → N$
  and $g : P × X → Y$ be
  smooth families of maps. If for each $p$ and for every $x$ not in $K$,
  $f_p(φ(x)) = ψ(g_p(x))$ then the family of maps $f_p$ updated using $g_p$ is
  smooth from $P × M$ to $N$.
\end{lemma}

\begin{proof}
  \leanok
  Note that $P × M = (P × φ(X)) ∪ (P × φ(K)^c)$. Both
  those sets are open and the updated maps coincide with
  $(p, m) ↦ ψ ∘ g_p ∘ φ⁻¹(m)$ on the first one and $f$ on the second one.
\end{proof}

\begin{lemma}
  \label{lem:dist_updating}
  \uses{def:update}
  \lean{OpenSmoothEmbedding.dist_update}
  \leanok
  Let $φ : X → M$ and $ψ : Y → N$ be smooth open embeddings.
  Let $K_X$ and $K_P$ be compact sets in $X$ and $P$.
  Let $f : P × M → N$ be a continuous family of
  maps such that, for each $p$, $f_p(φ(X)) ⊂ ψ(Y)$.
  For every continuous function $ε : M → ℝ_{> 0}$, there is some positive
  number $η$ such that, for every map $g : P × X → Y$ and every $(p, p', x)$ in
  $K_P × K_P × K_X$, $d(g_{p'}(x), ψ⁻¹∘f_p∘φ(x)) < η$ implies
  $d(f'_{p'}(φ(x)), f_p(φ(x))) < ε(φ(x))$ where $f'$ is obtained by updating $f$
  using $g$.
\end{lemma}

\begin{proof}
  \leanok
  Let $ε$ be a positive continuous function on $M$. Since $K_X$ is compact, we get
  a positive number $ε₀$ such that $ε(m) ≥ ε₀$ for each $m$ in $K_X$. We denote by
  $K₁$ the closed $1$-thickening of the image of $K_P × K_X$ under
  $(p, x) ↦ ψ⁻¹∘f_p∘φ(x)$. This is a compact set
  so $ψ$ is uniformly continuous on $K₁$ and we get a positive $τ$ such that
  for all $x$ and $y$ in $K₁$, $d(x, y) < τ ⇒ d(ψ(x), ψ(y)) < ε₀$.

  We now prove that $η = \min(τ, 1)$ is suitable. Fix $(p, p', x)$ in
  $K_P × K_P × K_X$ such that $d(g_{p'}(x), ψ⁻¹∘f_p∘φ(x)) < η$.
  In particular $d(g_{p'}(x), ψ⁻¹∘f_p∘φ(x)) < 1$ hence $g_{p'}(x)$ is in $K₁$. Since
  $ψ⁻¹∘f_p∘φ(x)$ is also in $K₁$ and $d(g_{p'}(x), ψ⁻¹∘f∘φ(x)) < τ$, we get
  $d(ψ∘ g_{p'}(x), ψ ∘ ψ⁻¹∘f_p∘φ(x)) < ε₀$. This precisely means that
  $d(f'_{p'}(φ(x)), f_p(φ(x)) < ε₀$. Since $(p, p', x)$ is in $K_P × K_P × K_X$,
  this is less than $ε(m)$.
\end{proof}

In order to handle in a uniform way compact and non-compact manifolds, we will
index sequences by the family of sets $\IT{N}$ defined for each natural number
$N$ by:
  \[
    \IT{N} =
    \begin{cases}
      ℕ \text{ if $N = 0$}\\
      \{0, \dots, N - 1\} \text{otherwise}
    \end{cases}
  \]

\begin{lemma}
  \label{lem:nice_atlas}
  \lean{nice_atlas}
  \uses{def:index_type}
  \leanok
  Let $M$ be a manifold modelled on the normed space $E$ and $(V_j)_{j ∈ J}$
  a cover of $M$ by open sets. There exists some natural number $N$ and
  a family of smooth open embeddings $φ : \IT{N} × E → M$ such that
  \begin{itemize}
    \item for each $i$ there is some $j$ such that $φ_i(E) \subseteq V_j$,
    \item $i ↦ φ_i(E)$ is a locally-finite collection of sets in $M$,
    \item $⋃_i φ_i(B_E(0, 1)) = M$ where $B_E(0, 1)$ is the open unit ball in $E$.
  \end{itemize}
\end{lemma}

\begin{proof}
  \leanok
  The proof is a standard compact-exhaustion argument.
  Let $K_0, K_1, K_2, \ldots $ be a compact exhaustion of $M$ and define:
  \begin{align*}
    C_n &= K_{n+2} \smallsetminus K_{n+1}^o,\\
    U_n &= K_{n+3}^o \smallsetminus K_n.
  \end{align*}
  Thus:
  \begin{itemize}
    \item $C_n$ is compact,
    \item $U_n$ is open,
    \item $C_n \subseteq U_n$,
    \item $\bigcup_n C_n = M$,
    \item $U_n \cap U_m = \emptyset$ if $|n - m| > 2$.
  \end{itemize}
  For any $y \in E$ and $r > 0$, fix a smooth diffeomorphism $f_{y,r} : E \simeq B_E(y, r)$ such
  that $f_{y,r}(0) = y$.
  For each $n$ and $x \in C_n$, let $\psi_x$  be a smooth chart mapping an open
  neighbourhood of $x$ to an open set of the model space $E$.
  Writing $y = \psi_x (x) \in E$, let:
  \begin{align*}
    B_{n, x} &= \psi_x^{-1} (B_E(y, r)),\\
    W_{n, x} &= \psi_x^{-1} (f_{y,r} (B_E (0, 1))),\\
  \end{align*}
  for some $r > 0$ (which may depend on $n$, $x$) sufficiently small that:
  \begin{itemize}
    \item $B_E(y, r)$ lies in the target of the chart $\psi_x$,
    \item $B_{n,x}$ is contained in $U_n$,
    \item $B_{n,x}$ is contained in $V_j$ for some $j$.
  \end{itemize}
  Note that $x \in W_{n,x}$. For each $n$, choose a finite subcovering of $C_n$ by
  $W_{n, x_1}, \ldots, W_{n, x_{l_n}}$ and define $\iota \subseteq ℕ \times M$ by:
  \begin{align*}
    \iota = \bigcup_n \{ (n, x_1), \ldots, (n, x_{l_n}) \} .
  \end{align*}
  Note that $\iota$ is countable and furthermore:
  \begin{itemize}
    \item for each $i \in \iota$, there is some $j$ such that $B_i \subseteq V_j$,
    \item $(B_i)_{i \in \iota}$ is locally-finite
    (indeed more is true: $B_i$ meets only finitely-many $B_{i'}$ for $i, i' \in \iota$
    since $B_{m, x} \cap B_{n, x'} = \emptyset$ if $|n - m| > 2$),
    \item $(W_i)_{i \in \iota}$ covers $M$.
  \end{itemize}
  Given $i = (n, x_j) \in \iota$, the required map $\phi_i : E \to M$ is just:
  \begin{align*}
    E \simeq B_E(y_j, r) \simeq B_{n,j} \subseteq M.
  \end{align*}
  Since $ι$ is countable, it is in bijection with some $\IT{N}$.
\end{proof}

\begin{definition}
  \label{def:localisation_data}
  \uses{def:index_type}
  \lean{LocalisationData}
  \leanok
  Let $f : M → N$ be a continuous map between manifolds. A localisation data
  for $f$ is a tuple $(E, F, N, ι, φ, ψ, j)$ where $E$ and $F$ are normed
  vector spaces, $N$ is a natural number, $ι$ is a set, $φ : \IT{N} × E → M$
  and $ψ : ι × F → N$  are families of smooth open embeddings, and
  $j : \IT{N} → ι$ such that:
  \begin{itemize}
    \item
      $⋃_i φ_i(B_E) = M$ where $B_E$ is the open unit ball in $E$,
      \item
      $⋃_i ψ_i(B_F) = N$ where $B_F$ is the open unit ball in $F$,
    \item
      $∀ i, f(φ_i(E)) ⊂ ψ_{j(i)}(B_F)$ where $B_F$ is the open unit ball in $F$,
    \item
      $i \mapsto ψ_i(F)$ is locally finite.
  \end{itemize}
  Such a tuple will be denoted by $(φ, ψ, j)$ for brevity.
\end{definition}

\begin{lemma}
  \label{lem:ex_localisation}
  \uses{def:localisation_data}
  \lean{stdLocalisationData}
  \leanok
  Any continuous map between manifolds has some localisation data.
\end{lemma}

\begin{proof}
  \uses{lem:nice_atlas}
  \leanok
  The preceding lemma (applied to the trivial cover of $N$ by itself) gives a
  family of $ψ : ι' × F → N$ of open smooth embeddings that the images of
  $B_F$ cover $N$. We then apply this lemma again to the cover of $M$ given by all
  $f⁻¹(ψ_j(B_F))$.
\end{proof}

The general idea will be to apply the results of the previous chapters to
all the $ψ_{j(i)}⁻¹ ∘ f ∘ φ_i : E → F$ for some maps $f$. However we must be
careful that doing this for some $i$ does not ruin the setup for the next $i$.
This is easier to control using a distance function on the target manifold as
in \Cref{lem:localisation_stability} below. First we need a general lemma about
a single metric space (actually the formalized statement is stronger, it assumes
only closed sets instead of compact ones, but here we explain the easier proof
which is sufficient for our purposes).

\begin{lemma}
  \label{lem:stability_cover}
  \lean{Metric.exists_continuous_real_forall_closedBall_subset}
  \leanok
  In a metric space $X$, let $U : ι → \set{X}$ be a
  family of open subsets of $X$ and let $K : ι → \set{X}$ be a locally-finite
  family of closed subsets such that $K_i ⊂ U_i$ for all $i$. There exists a continuous
  function $δ : X → ℝ_{> 0}$ such that:
  \[
    ∀ x\, x',\; ∀ i, \left[x ∈ K_i \text{ and } d(x, x') < δ(x)\right] ⇒ x' ∈ U_i.
  \]
\end{lemma}

\begin{proof}
  \leanok
  We first note that, for any given $i$, compactness of $K$ and openness of
  $V_i$ give a positive number $δ_i$ such that the $δ_i$-neighborhood of $K_i$
  is contained in $V_i$. We now prove that solutions exist locally. Let $x$ be
  any point in $X$. From the local finiteness assumption, we get a neighborhood
  $U$ of $x$ such that $\{i | U ∩ V_i ≠ ∅\}$ is finite. The constant function
  with value the minimum of the corresponding $δ_i$ is a solution on $U$. Since
  the condition we put on $δ$ is convex, we can glue those local solutions using
  \Cref{lem:exists_cont_diff_of_convex}.
\end{proof}

\begin{lemma}
  \label{lem:localisation_stability}
  \lean{localisation_stability}
  \leanok
  \uses{def:localisation_data}
  Let $f : M → N$ be a continuous map between manifolds, and let
  $(φ, ψ, i)$ be some localisation data for $f$. There exists a continuous
  positive function $ε : M → ℝ_{>0}$ such that:
  \[
    ∀ g : M → N, \big[∀ m,\; d(f(m), g(m)) < ε(m)\big] ⇒ ∀ i,\; g(φ_i(E)) ⊂ ψ_{j(i)}(F).
  \]
\end{lemma}
Note that, in the preceding lemma, the conclusion $g(φ_i(E)) ⊂ ψ_{j(i)}(F)$ is
weaker than the condition $f(φ_i(E)) ⊂ ψ_{j(i)}(B_F)$ that appears in the
definition of localisation data.

The condition $∀ m,\; d(g(m), f(m)) < ε(m)$ will be abbreviated $d(g, f) < ε$.

\begin{proof}
  \leanok
  \uses{lem:stability_cover}
  The preceding lemma applied to the family of open sets $ψ_j(F)$ and the
  family of compact sets $ψ_j(\overline{B_F})$ give a positive continuous
  function $δ : N → ℝ$ such that $ε = δ ∘ f$ is suitable. Indeed, assume
  $g : M → N$ satisfies $d(g, f) < ε$ and fix some $i$ and some $m ∈ φ_i(E)$.
  We know $f(m) ∈ ψ_{j(i)}(\overline{B_F})$ and our assumption on $g$ gives
  $d(g(m), f(m)) < δ(f(m))$. So the property of $δ$ ensures $g(m) ∈ ψ_{j(i)}(F)$.
\end{proof}


\subsection{Jets spaces}

We now need to introduce the bundles that will replace the jet spaces
$E × F × \Hom(E, F)$ from the previous chapter. We need a couple of
fiber bundles constructions.

\begin{definition}
\label{def:pull_back_bundle}
\leanok
\lean{Bundle.Pullback}
For every bundle $p : E → B$ and every map $f \co B' → B$,
the pull-back bundle $f^*E → B'$ is defined by
$f^*E = \{(b', e) ∈ B' × E \;|\; p(e) = f(b')\}$ with
the obvious projection to $B'$.
\end{definition}

\begin{definition}
\label{def:hom_bundle}
\leanok
\lean{Bundle.ContinuousLinearMap.vectorBundle}
Let $E → B$ and $F → B$ be two vector bundles over some smooth manifold
$B$. The bundle $\Hom(E, F) → B$ is the set of linear maps from
$E_b$ to $F_b$ for some $b$ in $B$, with the obvious projection map.
\end{definition}

Set-theoretically, one can define $\Hom(E, F)$ as the set of subsets
$S$ of $E × F$ such that there exists $b$ such that $S ⊂ E_b × F_b$
and $S$ is the graph of a linear map. But the type theory formalization
will use other tricks here. The facts that really matter are listed in
\cref{lem:one_jet_extension_prop} below.

\begin{definition}
\label{def:one_jet_space}
\leanok
\lean{OneJetBundle}
\uses{def:pull_back_bundle, def:hom_bundle}
Let $M$ and $N$ be smooth manifolds. Denote by
$p_1$ and $p_2$ the projections of $M × N$ to
$M$ and $N$ respectively.

The space $J^1(M, N)$ of $1$-jets of maps from $M$ to $N$ is
$Hom(p_1^*TM, p_2^*TN)$
\end{definition}

We will use notations like $(m, n, φ)$ to denote an element
of $J^1(M, N)$,
but one should keep in mind that $J^1(M, N)$ is not a product,
since $φ$ lives in $\Hom(T_mM, T_nN)$ which depends on $m$ and $n$.


\begin{definition}
\label{def:one_jet_extension}
\uses{def:one_jet_space}
\leanok
\lean{oneJetExt}
The $1$-jet of a smooth map $f \co M → N$ is the map from
$m$ to $J^1(M, N)$ defined by $j^1f(m) = (m, f(m), T_mf)$.
\end{definition}

The composition of a section $\F \co M → J^1(M, N)$ with the projection
onto $N$ will sometimes be denoted by $\bs \F \co M → N$ and called the
base map of $\F$. For any $m$, $\F(m)_φ$ will denote the component of $\F(m)$
living in $\Hom(T_mM, T_{\bs\F(m)}N)$.

\begin{lemma}
\label{lem:one_jet_extension_prop}
\lean{ContMDiff.oneJetExt}
\leanok
\uses{def:one_jet_extension}
For every smooth map $f \co M → N$,
\begin{enumerate}
  \item
    \label{lem:one_jet_smooth}\uses{def:one_jet_space, def:one_jet_extension}
    $j^1f$ is smooth
  \item
    \label{lem:one_jet_section}\uses{def:one_jet_space, def:one_jet_extension}
    $j^1f$ is a section of $J^1(M, N) → M$
\end{enumerate}
\end{lemma}

\begin{proof}
  \leanok
  Points 2 and 3 are obvious by construction.

  To show that $j^1f$ is smooth, suppose that $M$ is modelled over $E$ with charts
  $C_x : M \to E$ and coordinate change functions
  $C_{x,x'}=C_{x'}C_x^{-1} : E \to E$ and similarly let $C'_y$ be charts for $N$.
  By construction of the 1-jet bundle, we need to check that for each $x_0$ the map
  $$x\mapsto TC'_{f(x),f(x_0)}\circ T(C'_{f(x)}fC_x^{-1})\circ
    T_{C_{x_0}(x)}(C_{x_0,x}): M \to L(E,E)$$
  is smooth at $x_0$ (we occasionally omit the point where the tangent maps are taken).
  For $x$ close to $x_0$ the coordinate changes are smooth, so we can write
  \begin{align*}
  TC'_{f(x),f(x_0)}\circ T(C'_{f(x)}fC_x^{-1})\circ T(C_{x_0,x})
  &= T_{C_{x_0}(x)}(C'_{f(x),f(x_0)}C'_{f(x)}fC_x^{-1}C_{x_0,x})\\
  &= T_{C_{x_0}(x)}(C'_{f(x_0)}fC_{x_0})
  \end{align*}
  This is smooth since $C'_{f(x_0)}fC_{x_0}$ is smooth.
\end{proof}

\begin{definition}
\label{def:holonomic_section}
\leanok
\lean{OneJetSec.IsHolonomicAt, OneJetSec.isHolonomicAt_iff}
\uses{def:one_jet_space}
A section $\F$ of $J^1(M, N) → M$ is called holonomic if it is the
$1$--jet of its base map.
Equivalently, $\F$ is holonomic if there exists
$f \co M → N$ such that $\F = j^1f$, since such a map is
necessarily $\bs \F$.
\end{definition}

\section{First order differential relations}

\begin{definition}
  \label{def:rel}
  \leanok
  \lean{RelMfld}
  \uses{def:one_jet_space}
  A first order differential relation for maps from $M$ to $N$ is a
  subset $\Rel$ of $J^1(M, N)$.
\end{definition}

\begin{definition}
  \label{def:formal_sol}
  \leanok
  \lean{FormalSol}
  \uses{def:rel, def:one_jet_extension}
  A formal solution of a differential relation $\Rel ⊂ J^1(M, N)$ is a
  section of $J^1(M, N) → M$ taking values in $\Rel$.
\end{definition}


\begin{definition}
  \label{def:htpy_formal_sol}
  \leanok
  \lean{HtpyFormalSol}
  \uses{def:formal_sol}
  A homotopy of formal solutions of $\Rel$ is a smooth family of sections
  $\F : ℝ × M → J^1(M, N)$ such that each $m ↦ \F(t, m)$ is a formal solution.
\end{definition}

The next definition will be used in cases where $X$ and $Y$ are vector spaces,
in order to relate the global theory to the local one.

\begin{definition}
  \label{def:transfer_map}
  \leanok % TODO check onejetsec!
  \lean{OpenSmoothEmbedding.transfer, OneJetSec.localize, RelMfld.localize}
  \uses{def:one_jet_space}
  Given manifolds $M$, $X$, $N$ and $Y$ and smooth open embeddings $g : Y → N$
  and $h : X → M$ we get a transfer map $ψ_{g, h} : J^1(X, Y) → J^1(M, N)$
  defined by
  \[
    ψ_{g, h}(x, y, φ) = (h(x), g(y), T_yg ∘ φ ∘ (T_xh)⁻¹)
  \]
  and an operator on sections which sends $\F : M → J^1(M, N)$ to
  $Ψ_{g, h}\F : X → J^1(X, Y)$ defined when $\bs\F(h(X)) ⊂ g(Y)$ by
  \[
    Ψ_{g, h}\F(x) = (x, g⁻¹ ∘ \bs\F ∘ h(x), (T_{g⁻¹ ∘ \bs\F ∘ h(x)}g)⁻¹ ∘ \F(h(x))_φ ∘ T_xh).
  \]
  Given a relation $\Rel ⊂ J^1(M, N)$, the induced relation in $J^1(X, Y)$
  is $ψ_{g, h}⁻¹\Rel$.
\end{definition}

The following is a localization lemma needed to take advantage of all the work
from the previous chapter.

\begin{lemma}
  \label{lem:transfer}
  \leanok
  \lean{OneJetSec.localize, isHolonomicAt_localize_iff, OneJetSec.localize_mem_iff}
  \uses{def:transfer_map, def:holonomic_section, def:formal_sol}
  In the situation of the previous definition, given a section $\F : M → J^1(M, N)$:
  \begin{itemize}
    \item $Ψ_{g, h}(\F)$ is a smooth section of $J^1(X, Y)$.
    \item $\F$ is holonomic on $s ⊂ h(X) ∩ \bs\F⁻¹(g(Y))$ if and only if $Ψ_{g, h}(\F)$
      is holonomic on $h⁻¹(s)$.
    \item $\F$ is a formal solution of $\Rel$ on $h(X) ∩ \bs\F⁻¹(g(Y)$ if and only if $Ψ_{g, h}(\F)$
      is a formal solution of the induced relation $Ψ_{g, h}⁻¹\Rel$.
  \end{itemize}
\end{lemma}

\begin{proof}
  \leanok
  \uses{lem:one_jet_extension_prop}
  The first point is clear by composition. In order to prove the second point
  while keeping notations under control, we set
  $f(x) = g⁻¹ ∘ \bs\F ∘ h$. Using this notation
  $Ψ_{g, h}\F(x) = (x, f(x), (T_{f(x)}g)⁻¹ ∘ \F(h(x))_φ ∘ T_xh)$. We have
  \begin{align*}
    T_x f &= T_{\bs\F ∘ h(x)}(g⁻¹) ∘ T_{h(x)}\bs\F ∘ T_x h\\
          &= \left(T_{f(x)}g\right)⁻¹ ∘ T_{h(x)}\bs\F ∘ T_x h
  \end{align*}
  hence $Ψ_{g, h}\F$ is holonomic at $x$ if and only if
  $\left(T_{f(x)}g\right)⁻¹ ∘ \F(h(x))_φ ∘ T_xh = \left(T_{f(x)}g\right)⁻¹ ∘ T_{h(x)}\bs\F ∘ T_x h$
  and this is equivalent to $\F(h(x))_φ = T_{h(x)}\bs\F$ which is the holonomy condition for
  $\F$ at $h(x)$.

  The third point is a direct consequence of the easy formula $ψ_{g, h} ∘ Ψ_{g, h}(\F) = F ∘ h$.
\end{proof}

\begin{definition}
  \label{def:h-princ}
  \leanok
  \lean{RelMfld.SatisfiesHPrinciple, RelMfld.SatisfiesHPrincipleWith}
  \uses{def:htpy_formal_sol}
  A first order differential relation $\Rel ⊂ J^1(M, N)$ satisfies the
  $h$-principle if every formal solution of $\Rel$ is homotopic to a
  holonomic one.
  It satisfies the parametric $h$-principle if, for every manifold $P$
  and every closed set $C$ in $P × M$, every family $\F : P × M → J^1(M, N)$ of formal
  solutions which are holonomic for $(p, m)$ near $C$
  is homotopic to a family of holonomic ones relative to $C$.
\end{definition}


\subsection*{Parametricity for free}

In many cases, relative parametric $h$-principles can be deduced from relative
non-parametric ones with a larger source manifold.
Let $X$, $P$ and $Y$ be manifolds, with $P$ seen a parameter space.
Denote by $Ψ$ the map from $J^1(X × P, Y)$ to $J^1(X, Y)$ sending $(x, p, y, ψ)$ to
$(x, y, ψ ∘ ι_{x, p})$ where $ι_{x, p} : T_xX → T_xX × T_pP$ sends $v$ to $(v, 0)$.

To any family of sections $F_p : x ↦ (f_p(x), φ_{p, x})$ of $J^1(X, Y)$, we
associate the section $\bar F$ of $J^1(X × P, Y)$ sending $(x, p)$ to
$\bar F(x, p) := (f_p(x), φ_{p, x} ⊕ ∂f/∂p(x, p))$.

\begin{lemma}
  \label{lem:param_trick}
  \uses{def:holonomic_section, def:formal_sol}
  \leanok
  \lean{FamilyOneJetSec.isHolonomicAt_uncurry, FamilyOneJetSec.uncurry_mem_relativize}
  In the above setup, we have:
  \begin{itemize}
    \item
      $\bar F$ is holonomic at $(x, p)$ if and only if $F_p$ is holonomic
      at $x$.
    \item
      $F$ is a family of formal solutions of some $\Rel ⊂ J^1(X, Y)$ if and
      only if $\bar F$ is a formal solution of $\Rel^P := Ψ^{-1}(\Rel)$.
  \end{itemize}
\end{lemma}

\begin{proof}
  \leanok
  For the first part, the derivative of $\bar F$ is $∂f/∂x(x, p) ⊕ ∂f/∂p(x, p)$, which is equal to
  $\bar F_φ$ iff $∂f/∂x(x, p) = f_φ$.

  The second part follows from $Ψ ∘ \bar F(x,p)=F_p(x)$.
\end{proof}

\begin{lemma}
  \label{lem:param_for_free}
  \uses{def:h-princ}
  \leanok
  \lean{RelMfld.SatisfiesHPrinciple.satisfiesHPrincipleWith}
  Let $\Rel$ be a first order differential relation for maps from $M$ to
  $N$.
  If, for every manifold with boundary $P$, $\Rel^P$ satisfies the
  $h$-principle then $\Rel$ satisfies the parametric $h$-principle.
  Likewise, the $C^0$-dense and relative $h$-principle for all
  $\Rel^P$ imply the parametric $C^0$-dense and relative $h$-principle for
  $\Rel$.
\end{lemma}

\begin{proof}
  \leanok
  \uses{lem:param_trick}
  By \cref{lem:param_trick} we can turn a formal solution of $\Rel$ into a formal solution of
  $\Rel^P$, so we get a homotopy to a holonomic formal solution. We can turn this homotopy back to
  a homotopy of the original formal solution.
\end{proof}


\section{The $h$-principle for open and ample differential relations}
\label{sec:general_theory}

In this chapter, $X$ and $Y$ are smooth manifolds and $\Rel$ is a first order
differential relation on maps from $X$ to $Y$: $\Rel ⊂ J^1(X, Y)$.
For any $σ = (x, y, φ)$ in $\Rel$ and any dual pair
$(λ, v) ∈ T^*_xX × T_xX$,
we set:
\[
    \Rel_{σ, λ, v} =
     \Conn_{φ(v)}\left\{w ∈ T_yY \;;\;
       \big(x,\; y,\; φ + (w - φ(v))⊗λ\big) ∈ \Rel\right\}
\]
where $\Conn_a A$ is the connected component of $A$ containing $a$. In order to
decipher this definition, it suffices to notice that $φ + (w - φ(v))⊗λ$ is the
unique linear map from $T_xX$ to $T_yY$ which coincides with $φ$ on $\ker λ$
and sends $v$ to $w$. In particular, $w = φ(v)$ gives back $φ$.

Of course we will want to deal with more than one point, so we will consider a
vector field $V$ and a $1$--form $λ$ such that $λ(V) = 1$ on some subset $U$ of
$X$, a formal solution $F$ (defined at least on $U$), and get the corresponding
$\Rel_{F, λ, v}$ over $U$.

One easily checks that $\Rel_{σ, κ^{-1}λ, κv} = κ\Rel_{σ, λ, v}$ hence the above
definition only depends on $\ker λ$ and the direction $ℝV$.

\begin{definition}
  \label{def:ample_relation}
  \uses{def:one_jet_space, def:ample_subset}
  \leanok
  \lean{RelMfld.Ample}
  A relation $\Rel$ is ample if, for every $σ = (x, y, φ)$ in $\Rel$ and every
  $(λ, v)$, the slice $\Rel_{σ, λ, v}$ is ample in $T_yY$.
\end{definition}

\begin{lemma}
  \label{lem:ample_iff_loc}
  \uses{def:ample_relation, def:transfer_map}
  \leanok
  \lean{RelMfld.Ample.localize}
  Given manifolds $W$, $X$, $Y$ and $Z$ and smooth open embeddings $g : Z → Y$
  and $h : W → X$, the relation induced (in the sense of \Cref{def:transfer_map})
  in $J^1(W, Z)$ by an ample relation in $J^1(X, Y)$ is ample.
\end{lemma}

\begin{proof}
  \leanok
  By definition, the relation induced by $\Rel$ is
  $ψ_{g, h}⁻¹\Rel$ where
  $ψ_{g, h}(w, z, φ) = (h(w), g(z), T_zg ∘ φ ∘ (T_wh)⁻¹)$.
  Fix $σ =(w, z, φ) ∈ ψ_{g, h}⁻¹\Rel$
  and a dual pair $(λ, v)$ on $T_wW$. We set $G = T_z g$ and $H = T_w h$. Both
  those maps are linear isomorphisms. We compute the slice corresponding to $(σ, λ, v)$:
  \begin{align*}
    ψ_{g, h}⁻¹\Rel(σ, λ, v)
    &= \Conn_{φv}\left\{u ∈ T_wW \;\left|\; (w, z, φ + (u - φv)⊗λ) ∈ ψ_{g, h}⁻¹\Rel \right.\right\}\\
    &= \Conn_{φv}\left\{u ∈ T_wW \;\left|\; (h(w), g(z), G ∘ (φ + (u - φv)⊗λ)∘ H⁻¹) ∈ \Rel \right.\right\}\\
%    &= \left\{u ∈ T_wW \;\left|\; (h(w), g(z), G ∘ φ ∘ H⁻¹ + (Gu - G ∘ φ ∘ H⁻¹ ∘ Hv)⊗(λ ∘ H⁻¹)) ∈ \Rel \right.\right\}\\
%    &= \left\{u ∈ T_wW \;\left|\; (h(w), g(z), G ∘ φ ∘ H⁻¹ + (Gu - G ∘ φ ∘ H⁻¹ ∘ Hv)⊗(λ ∘ H⁻¹)) ∈ \Rel \right.\right\}\\
%    &= \left\{u ∈ T_wW \;\left|\; Gu ∈ \Rel((h(w), g(z), G ∘ φ ∘ H⁻¹), λ ∘ H⁻¹, Gu) \right.\right\}\\
    &= G⁻¹\Rel(ψ_{g, h}σ, λ ∘ H⁻¹, Hv).
  \end{align*}
  Hence the slice $ψ_{g, h}⁻¹\Rel(σ, λ, v)$ is the image of a slice of $\Rel$
  under a linear isomorphism, hence ample.
\end{proof}

\begin{lemma}
  \label{lem:open_ample_immersion}
  \uses{def:ample_relation}
  \leanok
  \lean{immersionRel_open_ample}
  The relation of immersions of $M$ into $N$ in positive codimension is open
  and ample.
\end{lemma}

\begin{proof}
  \proves{lem:open_ample_immersion}
  \uses{lem:ample_codim_two}
  \lean{immersionRel_open_ample}
  \leanok
  For every $σ = (x, y, φ)$ in the immersion relation $\Rel$,
  and for every dual pair $(π, v)$, the slice
  $\Rel(σ, π, v)$ is the set of $w$ which do not belong to
  the image of $\ker π$ under $φ$.
  Since $\dim M > \dim N$, this image has codimension at least $2$ in
  $T_yN$, and \cref{lem:ample_codim_two} concludes.
\end{proof}



We first explain how to get rid of parameters, using the
relation $\Rel^P$ for families of solutions parametrized by $P$.

\begin{lemma}
    \label{lem:ample_parameter}
    \uses{def:ample_relation}
    \lean{RelMfld.Ample.relativize}
    \leanok
    If $\Rel$ is ample then, for any parameter space $P$, $\Rel^P$ is also ample.
\end{lemma}

\begin{proof}
  \leanok
  We fix $σ = (x, y, ψ)$ in $\Rel^P$.
  For any $λ = (λ_X, λ_P) ∈ T^*_xX × T^*_pP$ and $v = (v_X, v_P) ∈ T_xX × T_pP$
  such that $λ(v) = 1$, we need to prove that $\Conv\Rel^P_{σ, λ, v} = T_yY$.
  Unfolding the definitions gives:
  \[
  \Rel^P_{σ, λ, v} = \Conn_{φ(v)}\left\{w ∈ T_yY \;;\;
      \big(x,\; y,\; ψ ∘ ι_{x, p} + (w - ψ(v))⊗λ_X\big) ∈ \Rel\right\}.
  \]
  A degenerate but easy case is when $λ_X = 0$. Then the condition on $w$
  becomes $ψ ∘ ι_{x, p} ∈ \Rel$, which is true by definition of $\Rel^P$, so
  $\Rel^P_{σ, λ, v} = T_yY$.

  We now assume $λ_X$ is not zero and choose $u ∈ T_xX$ such that $λ_X(u) = 1$.
  We then have $\Rel^P_{σ, λ, v} = \Rel_{Ψσ, λ_X, u} + ψ(v) - ψ∘ι_{x, p}(u)$.
  Because $\Rel$ is ample and taking convex hull commutes with translation, we
  get that $\Conv\Rel^P_{σ, λ, v} = T_yY$.
\end{proof}



\begin{theorem}[Smale 1958]
  \label{thm:sphere_eversion}
  \lean{sphere_eversion}\leanok
	There is a homotopy of immersions of $𝕊^2$ into $ℝ^3$ from the inclusion map to
	the antipodal map $a : q ↦ -q$.
\end{theorem}

\begin{proof}
  \leanok
  \uses{thm:open_ample, lem:open_ample_immersion}
	We denote by $ι$ the inclusion of $𝕊^2$ into $ℝ^3$.
	We set $j_t = (1-t)ι	+ ta$.
  This is a homotopy from $ι$ to $a$ (but not an immersion for $t=1/2$).
  Using the canonical trivialization of the tangent
	bundle of $ℝ^3$, we can set, for $(q, v) ∈ T𝕊^2$,
	$G_t(q, v) = \mathrm{Rot}_{Oq}^{πt}(v)$, the rotation around axis $Oq$ with
	angle $πt$.
  The family $σ : t ↦ (j_t, G_t)$ is a homotopy of formal immersions
  relating $j^1ι$ to $j^1a$.
  It is homotopic by reparametrization to a homotopy of formal immersions
  relating $j^1ι$ to $j^1a$ which are holonomic for $t$ near the
  $0$ and $1$.

  The above theorem ensures this family is homotopic,
	relative to $t = 0$ and $t = 1$, to a family of holonomic formal immersions,
	ie a family $t ↦ j^1f_t$ with $f_0 = ι$, $f_1 = a$, and each $f_t$ is an
	immersion.
\end{proof}

% vim: set expandtab sw=2 tabstop=2 tw=80:

\chapter{Local sphere eversion}
\label{chap:local_eversion}

The local theory of \Cref{chap:local} is already enough to deduce
Smale's sphere eversion theorem, although it is less natural than going through
the general results of \Cref{chap:global}. The goal of this appendix is to
explain how to do so. In this section $E$ denote a finite dimensional real
vector space equipped with an inner product. Later we will assume it is
3-dimensional. We denote by $𝕊$ the unit sphere in $E$.

Although we want to study immersions of $𝕊$ into $E$, we want to work only with
functions defined on the whole $E$. So we introduce a slightly artificial relation.
We denote by $B$ the open ball with radius $9/10$ around the origin in $E$ and set:

\[
  \Rel := \{(x, y, φ) ∈ J^1(E, E) \;|\; x ∉ B ⇒ \rst{φ}{x^⊥} \text{ is injective}\}.
\]

Of course solutions of this relation restrict to immersions of $𝕊$.

\begin{lemma}
  \label{lem:loc_immersion_rel_open}
  \lean{loc_immersion_rel_open}
  \leanok
  The relation $\Rel$ above is open.
\end{lemma}

\begin{proof}
  \leanok
  The main task is to fix $x_0 \notin B$
  and $\varphi_0 \in L(E, E)$ which is injective on $x_0^\perp$ and prove that,
  for every $x$ close to $x_0$ and $\varphi$ close to $\varphi_0$, $\varphi$ is
  injective on $x^\perp$. This is a typical situation where geometric intuition
  makes it feel like there is nothing to prove.

  One difficulty is that the subspace $x^\perp$ moves with $x$. We reduce to a fixed
  subspace by considering the restriction to $x_0^\perp$ of the orthogonal
  projection onto $x^\perp$. One can check this is an isomorphism as long as $x$
  is not perpendicular to $x_0$.
  More precisely, we consider $f \co J^1(E, E) \to ℝ \times L(x_0^\perp, E)$ which sends
  $(x, y, \varphi)$ to $(\langle x_0, x\rangle, \varphi \circ \pr_{x^\perp} \circ j_0)$
  where $j_0$ is the inclusion of $x_0^\perp$ into $E$. The set $U$ of injective
  linear maps is open in $L(x_0^\perp, E)$ and the map $f$ is continuous
  hence the preimage of $\{0\}^c \times U$ is open. This is good enough for us because
  injectivity of $\varphi \circ \pr_{x^\perp} \circ j_0$ implies injectivity of
  $\varphi$ on the image of $\pr_{x^\perp} \circ j_0$ which is $x^\perp$ whenever
  $\langle  x_0, x\rangle \neq 0$.
\end{proof}

\begin{lemma}
  \label{lem:loc_immersion_rel_ample}
  \lean{loc_immersion_rel_ample}
  \uses{lem:ample_codim_two}
  \leanok
  The relation $\Rel$ above is ample.
\end{lemma}

\begin{proof}
  \leanok
  The core fact here is that if one fixes vector spaces $F$ and $F'$, a dual
  pair $(\pi, v)$ on $F$ and an injective linear map $\varphi \co F \to F'$
  then the updated map $\Upd{p}{\varphi}{w}$ is injective if and only if $w$
  is not in $\varphi(\ker\pi)$. First we assume $\Upd{p}{\varphi}{\varphi(u)}$
  is injective for some $u$ in $\varphi(\ker\pi)$ and derive a contradiction.
  We have $\Upd{p}{\varphi}{\varphi(u)}v = \varphi(u)$ by the general
  definition of updating and also $\Upd{p}{\varphi}{\varphi(u)}u = \varphi(u)$
  since $u$ is in $\ker \pi$. Hence injectivity of $\varphi$ ensure $u = v$,
  which is absurd since $\pi(u) = 0$ and $\pi(v) = 1$. Conversely assume $w$
  is not in $\varphi(\ker\pi)$ and let us prove $\Upd{p}{\varphi}{w}$ is
  injective. Assume $x$ is in the kernel of $\Upd{p}{\varphi}{w}$. Decompose
  $x = u + tv$ with $u \in \ker\pi$ and $t$ a real number. We have
  $\Upd{p}{\varphi}{w}(x) = \varphi(u) + tw$. Hence our assumption on $x$
  implies $t$ vanishes otherwise we would have $w = -t⁻¹\varphi(u)$
  contradicting that $w$ isn't in $\varphi(\ker\pi)$. This vanishing and the
  assumption on $x$ then imply $\varphi(u) = 0$. Since $\varphi$ is injective
  we conclude that $u = 0$ and finally $x = 0$.

  We now turn to $\Rel$. It suffices to prove that for every $\sigma = (x, y,
  \varphi) \in \Rel$ and every dual pair $p = (\pi, v)$ on $E$, the slice
  $\Rel(\sigma, p)$ is ample. If $x$ is in $B$ then $\Rel(\sigma, p)$ is the
  whole $E$ which is obviously ample. So we assume $x$ is not in $B$. Since
  $\sigma$ is in $\Rel$, $\varphi$ is injective on $x^\perp$. The slice is the
  set of $w$ such that $\Upd{p}{\varphi}{w}$ is injective on $x^\perp$. Assume
  first $\ker\pi = x^\perp$. Then $\Upd{p}{\varphi}{w}$ coincides with
  $\varphi$ on $x^\perp$ hence the slice is the whole $E$. Assume now that
  $\ker\pi \neq x^\perp$. The slice is not very easy to picture in this case.
  But one should remember that, up to affine isomorphism, the slice depends
  only on $\ker \pi$. More specifically, if we keep $\pi$ but change $v$ then
  the slice is simply translated in $E$. Here we replace $v$ by the projection
  on $x^\perp$ of the vector dual to $\pi$ rescaled to keep the property
  $\pi(v) = 1$. What has been gained is that we now have $v \in x^\perp$ and
  $x^\perp = (x^\perp \cap \ker \pi) \oplus ℝ v$. Since $\varphi$ is
  injective on $x^\perp$, $\varphi(x^\perp \cap \ker \pi)$ is a hyperplane
  in $x^\perp$ and $\Upd{p}{\varphi}{w}$ is injective on $x^\perp$ if and only
  if $w$ is in the complement of $\varphi(x^\perp \cap \ker \pi)$ according to
  the core fact above. Since it is an hyperplane in $x^\perp$, it has
  codimension at least $2$ in $E$ hence its complement is ample.
\end{proof}

\begin{theorem}[Smale 1958]
  \lean{sphere_eversion_of_loc}
  \label{sphere_eversion_of_loc}
  \leanok
  There is a homotopy of immersion of $𝕊^2$ into $ℝ^3$ from the inclusion map to
	the antipodal map $a \co q ↦ -q$.
\end{theorem}

\begin{proof}
  \leanok
  \uses{lem:loc_immersion_rel_open, lem:loc_immersion_rel_ample, lem:h_principle_open_ample_loc, lem:param_trick}
	We denote by $ι$ the inclusion of $𝕊^2$ into $ℝ^3$.
	We set $j_t = (1-t)ι	+ ta$.
  This is a homotopy from $ι$ to $a$ (but not an immersion for $t=1/2$).
  Using the canonical trivialization of the tangent
	bundle of $ℝ^3$, we can set, for $(q, v) ∈ T𝕊^2$,
	$G_t(q, v) = \mathrm{Rot}_{Oq}^{πt}(v)$, the rotation around axis $Oq$ with
	angle $πt$.
  The family $σ : t ↦ (j_t, G_t)$ is a homotopy of formal immersions
  relating $j^1ι$ to $j^1a$. Those formal solutions are holonomic when $t$ is
  zero or one, so we can reparametrize the family to make such it is holonomic
  when $t$ is close to zero or one. Then we can extend it to a homotopy of
  formal solutions of $\Rel$ using a suitable cut-off ensuring smoothness
  near the origin. The relation $\Rel$ is ample according to
  \Cref{lem:loc_immersion_rel_ample} and then \Cref{lem:param_trick} ensures
  its 1-parameter version $\Rel^ℝ$ is also ample. The relation $\Rel$ is open according to
  \Cref{lem:loc_immersion_rel_open} hence $\Rel^ℝ$ is also ample.
  So we can use \Cref{lem:h_principle_open_ample_loc} to deform our family of
  formal solutions into a holonomic one.
\end{proof}

% vim: set expandtab sw=2 tabstop=2:

\chapter{From local to global}%
\label{cha:from_local_to_global}

In this chapter, we gather some topological preliminaries allowing to build
global objects from local ones. This is usually not discussed in informal expositions
where such arguments are either implicit or interspersed with more specific arguments.

We first need to discuss how to build a function having everywhere some local
properties from a sequence of functions having those properties on bigger and
bigger parts of the source space. We actually want to also accommodate finite sequences
so we start with a definition of the source of our sequences.

\begin{definition}
  \label{def:index_type}\lean{IndexType}\leanok
  For every natural number $N$ we set
  \[
    \IT{N} =
    \begin{cases}
      ℕ \text{ if $N = 0$}\\
      \{0, \dots, N - 1\} \text{otherwise}
    \end{cases}
  \]
\end{definition}

On each $\IT{N}$ we use the obvious linear ordering. In particular there is no
maximal element when $N = 0$ and $N-1$ is maximal if $N$ is positive.
The successor function $S \co \IT{N} → \IT{N}$ is the function sending $n$ to
$n+1$ unless $n$ is maximal, in which case $S(n) = n$.

Our first lemma gives a criterion ensuring that a sequence of functions is locally
ultimately constant hence has a limit that locally ultimately agrees with the
elements of the sequence.
Remember that a family of sets $V_n$ in a topological space $X$ is locally finite if
every point of $X$ has a neighborhood that intersects only finitely many $V_n$.

\begin{lemma}
  \label{lem:exists_forall_eventually_of_index_type}
  \uses{def:index_type}
  \lean{LocallyFinite.exists_forall_eventually_of_indexType}
  \leanok
  Let $X$ be a topological space and let $Y$ be any set. Let
  $f$ be a sequence of functions from $X$ to $Y$ indexed by $\IT{N}$
  for some $N$. Let $V$ be a family of subsets of $X$ indexed by $\IT{N}$
  such that, for every non-maximal $n$, $f_{S(n)}$ coincides with $f_n$ outside
  $V_{S(n)}$. If $V$ is locally finite then there exists $F \co X → Y$ such
  that, for every $x$ and every sufficiently large $n$,
  $F$ coincides with $f_n$ near $x$.
\end{lemma}

\begin{proof}\leanok
  The assumption that $V$ is locally finite gives, for every $x$ in $X$, a
  subset $U_x$ of $X$ such that $U_x$ is a neighborhood of $x$ and intersects
  only finite many $V_n$'s. In particular we can find an upper bound $n₀(x)$ of the set
  of indices $n$ in $\IT{N}$ such that $V_n$ intersects $U_X$.
  Since, for every non-maximal $n$, $f_{S(n)}$ coincides with $f_n$ outside
  $V_{S(n)}$, we get by induction that, for all $n ≥ n₀(x)$, $f_n$ coincides
  with $f_{n₀(x)}$ on $U_x$.

  We now define $F$ as $x ↦ f_{n₀(x)}(x)$. We claim that, for every $x$,
  $F$ coincides with $f_n$ on $U_x$ as soon as $n$ is at least $n₀(x)$.
  Indeed let us fix $x$ and $n ≥ n₀(x)$ and $y ∈ U_x$. We have
  \begin{align*}
    f_n(y) &= f_{n₀(x)}(y) \text{ since $n ≥ n₀(x)$ and $y ∈ U_x$}\\
           &= f_{\max(n₀(x), n₀(y))}(y) \text{ since $\max(n₀(x), n₀(y)) ≥ n₀(x)$ and $y ∈ U_x$}\\
           &= f_{n₀(y)}(y) \text{ since $\max(n₀(x), n₀(y)) ≥ n₀(y)$ and $y ∈ U_y$}\\
           &= F(y) \text{ by definition of $F$}.
  \end{align*}
\end{proof}

In the preceding lemma, the limit function $F$ inherits all local properties of
the elements of the sequence. In order to make this precise, we need the
language of germs of functions. One can define germs with respect to any filter but
we will need only the case of neighborhood filters~: two functions $f$ and $g$
define the same germ at some point $x$ if they coincide near $x$.

\begin{definition}\label{def:germ}\lean{Filter.Germ}\leanok
  Let $X$ be a topological space, $x$ a point in $X$ and $Y$ a set. A germ of function
  from $X$ to $Y$ at $x$ is an element of the quotient $(X → Y)_x$ of the set of functions
  from $X$ to $Y$ by the relation $f ∼ g$ if $f$ and $g$ coincide near $x$. The image
  of a function $f$ in this quotient will be denoted by $[f]_x$.

  A local predicate on functions from $X$ to $Y$ is a family $P$ of predicates on the
  germ set $(X → Y)_x$ for every $x$ in $X$. We say that a function $f$
  satisfies $P$ at $x$ if $P [f]_x$ holds, and $f$ satisfies $P$ if it
  satisfies $P$ at every point.
\end{definition}

For instance if $Y$ is also equipped with a topology then continuity is
(equivalent to) a local predicate on functions from $X$ to $Y$ since a function
is continuous if and only if it is continuous at every point $x$ and this
condition only depends on the germ of the function at $x$.

We also need to build local predicates by localizing some local predicates near
some subsets.

\begin{definition}
  \label{def:restrict_germ_predicate}\lean{RestrictGermPredicate}\leanok
  Let $X$ be a topological space, $A$ a subset of $X$, $Y$ a set and $P$ a
  local predicate on functions from $X$ to $Y$. The restriction of $P$ to $A$
  is the local predicate $P_{|A}$ defined by the constraint that a function
  $f$ satisfies $P_{|A}$ at $x$ if $x ∈ A$ implies that $f$ satisfies $P$ near $x$.
\end{definition}

Note the above definition hides a little lemma asserting that the obtained
predicate is indeed local. An even smaller lemma asserts that a function satisfies
$P_{|A}$ if and only if it satisfies $P$ at each point near $A$.

In the next lemma, there are three predicates or families of predicates. The
local predicate $P₀$ is satisfied by every function appearing in the lemma, it
could be a continuity or smoothness constraint. The family of local predicates
$P₁$ is the main constraint and the goal is to build a function satisfying all
of them. The family of predicates $P₂$ plays an auxiliary role, it does not
have to be local, does not appear in the conclusion and is only used to bring more
flexibility in the main inductive assumption. One can read ``$f$ satisfies
$P₂^i$'' as ``$f$ can be improved in $Uᵢ$''.

\begin{lemma}
  \uses{def:index_type}
  \label{lem:inductive_construction}\uses{def:germ}\lean{inductive_construction}\leanok
  Let $X$ be a topological space and $Y$ be any set. Let $U$ be a locally
  finite family of subsets of $X$ indexed by some $\IT{N}$. Let $P₀$ be a local
  predicate on functions from $X$ to $Y$, let $i ↦ P₁^i$ be a family of such
  predicates, and let $i ↦ P₂^i$ be a family of predicates on functions from $X$ to
  $Y$, all families being indexed by $\IT{N}$. Assume that
  \begin{itemize}
    \item
      there exists $f₀ \co X → Y$ satisfying $P₀$ and $P₂^0$~;
    \item
      for every $i$ in $\IT{N}$ and every $f \co X → Y$ satisfying $P₀$, $P₂^i$
      and every $P₁^j$ for $j < i$, there exists a function $f' \co X → Y$
      which coincides with $f$ outside $U_i$ and satisfies $P₀$ and every
      $P₁^j$ for $j ≤ i$ as well as $P₂^{S(i)}$ unless $i$ is maximal.
  \end{itemize}
  Then there exists $f \co X → Y$ which satisfies $P₀$ and all $P₁^i$'s.
\end{lemma}

\begin{proof}
  \leanok\uses{lem:exists_forall_eventually_of_index_type}
  The main assumption from the lemma allows to build by induction a sequence
  $f$ of functions from $X$ to $Y$ indexed by $\IT{N}$ such that, for every $n ∈ \IT{N}$,
  \begin{itemize}
    \item $f_n$ satisfies $P₀$
    \item for every $i ≤ n$, $f_n$ satisfies $P₁^i$.
    \item $f_{S(n)}$ satisfies $P₂$ unless $n$ is maximal.
    \item $f_{S(n)}$ coincides with $f_n$ outside $U_{S(n)}$.
  \end{itemize}
  Note that the first term of this sequence isn't $f₀$ but the function obtained by
  applying the induction assumption to $f₀$.

  The preceding lemma applied to this sequence gives a map $f$ which locally
  coincides with every element which is far enough in the sequence.
  Let $x$ be a point in $X$. Let $n$ be large enough to ensure $f$ coincides with
  $f_n$ near $x$. By definition this means $[f]_x = [f_n]_x$ and we know $P₀ [f_n]_x$
  hence we get $P₀ [f]_x$. Now fix also $n$ in $\IT{N}$. Let $n'$ be large
  enough to be larger than $n$ and such that $[f]_x = [f_{n'}]_x$. Since $n' ≥ n$ we have
  $P₁^n [f_{n'}]_x$ hence $P₁^n [f]_x$.
\end{proof}

Next we will need a version of the above lemma building a homotopy of maps.
In this version, $P₀$ is still a predicate such as continuity satisfied by
all maps from $X$ to $Y$ entering the discussion. Then $P₀'$ is analogous but for
maps from $ℝ × X$ to $Y$, and it will come with some affine invariance assumption
ensuring its compatibility with concatenation of homotopies.
Instead of having a completely general family of local predicates
$P₁^i$, we fix a single local predicate $P₁$ but it will be required to hold only
near some subset $K_i$ (as in \Cref{def:restrict_germ_predicate}).

Homotopies of maps from $X$ to $Y$ are usually meant to be continuous maps from
$[0, 1] × X$ to $Y$. In a differential topology context, one requires
smoothness and in order to be able to easily concatenate homotopies, it is very
convenient to add the assumptions that those maps are independent of the time
variable $t ∈ [0, 1]$ when $t$ is close to $0$ or $1$. Especially in a
formalization context, it is even more convenient to assume homotopies are defined on
$ℝ × X$, and time independent near $(-∞, 0] × X$ and $[1, +∞) × X$. Continuity or
smoothness don't appear in the following abstract lemma where they are replaced by
arbitrary local predicates.

\begin{lemma}
  \label{lem:inductive_htpy_construction}\lean{inductive_htpy_construction}\leanok
  Let $X$ be a topological space and $Y$ be any set.
  Let $P₀$ and $P₁$ be local predicates on maps from $X$ to $Y$.
  Let $P₀'$ be a local predicate on maps from $ℝ × X → Y$. Assume that
  for every $a$, $b$ and $t$ in $ℝ$, every $x$ in $X$ and every
  $f \co ℝ × X → Y$, if $f$ satisfies $P₂$ at $(at + b, x)$ then
  $(t, x) ↦ f(at+b, x)$ satisfies $P₀'$ at $(t, x)$. Let $f₀ \co X → Y$ be a function
  satisfying $P₀$ and such that $(t, x) ↦ f₀(x)$ satisfies $P₀'$.

  Let $K$ and $U$ be families of subsets of $X$ indexed by some $\IT{N}$.
  Assume that $U$ is locally finite and $K$ covers $X$.

  Assume that, for every $i$ in $\IT{N}$ and every $f \co X → Y$ satisfying $P₀$ and
  satisfying $P₁$ on $\bigcup_{j < i} K_j$, there exists $F \co ℝ × X → Y$ such that
  \begin{itemize}
    \item for all $t$, $F(t, \cdot)$ satisfies $P₀$
    \item $F$ satisfies $P₀'$
    \item $F(1, \cdot)$ satisfies $P₁$ on $\bigcup_{j ≤ i} K_j$
    \item $F(t, x) = f(x)$ whenever $x$ is not in $U_i$ or $t$ is near $(-∞, 0]$
    \item $F(t, x) = F(1, x)$ whenever $t$ is near $[1, +∞)$.
  \end{itemize}
  Then there exists $F \co ℝ × X → Y$ such that
  \begin{itemize}
    \item for all $t$, $F(t, \cdot)$ satisfies $P₀$
    \item $F$ satisfies $P₀'$
    \item $F(0, \cdot) = f₀$
    \item $F(1, \cdot)$ satisfies $P₁$.
  \end{itemize}
\end{lemma}

\begin{proof}\leanok\uses{lem:inductive_construction, def:restrict_germ_predicate}
  Carefully checking all details is a bit technical but the strategy is as follows.
  We fix an increasing sequence $T \co \IT{N} → [0, 1)$ starting at $0$,
  say $i ↦ 1 - 1/2^i$.
  We want to build a sequence of homotopies $F_i \co ℝ × X → Y$ where each
  $F_i$ is time-independent on $[T_i, + ∞) × X$ and, assuming $i$ isn't
  maximal, $F_{S(i)}$ is built from $F_i$ by applying the induction assumption
  to $F_i(T_i, \cdot)$ and rescaling the obtained homotopy by the affine map
  sending $[0, 1]$ to $[T_i, T_{S(i)}]$.

  Hence we want to apply \Cref{lem:inductive_construction} with source space
  $\hat{X} = ℝ × X$.
  We use as the background local predicate $\hat{P}₀$ at $(t, x)$ the
  constraint on a function $F$ that $F(t, \cdot)$ satisfies $P₀$ at $x$ and $P₀'$ at
  $(t, x)$ and if $t = 0$ then $F(t, x) = f₀(x)$. As the target family of local predicates
  $\hat{P}₁$ we use for every $i ∈ \IT{N}$ the constraint on $F$ at $(t, x)$
  that if $t = 1$ and $x$ is near $K_i$ then $F(t, \cdot)$ should satisfy $P₂$
  at $x$. As the auxiliary family of predicates $\hat{P₂}$ at index $i$ we use
  the constraint of being time-independent on $[T_i, + ∞) × X$.

  In order to explain how the induction assumption of the current lemma implies
  the induction assumption of \Cref{lem:inductive_construction}, we fix $i ∈ \IT{N}$
  and a map $F \co ℝ × X → Y$ that satisfies $\hat{P}₀$, is time-independent on
  $[T_i, + ∞) × X$ and satisfies $\hat{P}₁^j$ for all $j < i$. By this last
  requirement and the time independence property, we get that $F(T_i, \cdot)$
  satisfies $P₁$ near $\bigcup_{j < i} K_j$. Our induction assumption applied to
  $F(T_i, \cdot)$ then gives $F' \co ℝ × X → Y$ such that
  \begin{itemize}
    \item for all $t$, $F'(t, \cdot)$ satisfies $P₀$
    \item $F'$ satisfies $P₀'$
    \item $F'(1, \cdot)$ satisfies $P₁$ on $\bigcup_{j ≤ i} K_j$
    \item $F'(t, x) = f(x)$ whenever $x$ is not in $U_i$ or $t$ is near $(-∞, 0]$
    \item $F'(t, x) = F(1, x)$ whenever $t$ is near $[1, +∞)$.
  \end{itemize}
  As the new map required by the inductive assumption of \Cref{lem:inductive_construction},
  we pick
  \[
    F'' \co (t, x) ↦
    \begin{cases}
      F(t, x) \text{ if $t ≤ T_i$}\\
      F'\left((t - T_i)/(T_{S(i)} - T_i), x\right) \text{ if $t > T_i$}
    \end{cases}
  \]

  Fully checking that $F''$ is suitable is fairly technical but mostly straightforward.
  Care is required in particular to check that $F''$ coincides with $F$ near
  $(T_i, x)$ for every $x$. This uses both the fact that $F$ is time-independent on
  $[T_i, + ∞) × X$ and that $F'$ is time-independent near $(-∞, 0] × X$ hence in particular
  near $(0, x)$.
\end{proof}

In a different direction, we need a version of \Cref{lem:inductive_construction} where
we do not fix any family of subsets of the source space, but simply want to derive existence
of a function satisfying some local predicates from the assumptions of
existence of local solution and the ability to patch solutions. This requires putting a lot
more constraints on the source topological space in order to use the following classical
result.

\begin{lemma}
  \label{lem:exists_locally_finite_subcover_of_locally}
  \lean{exists_locallyFinite_subcover_of_locally}\leanok
  Let $X$ be a metrizable locally compact second countable topological space.
  Let $C$ be a closed subset in $X$. Let $P$ be a non-decreasing predicate on subsets
  of $X$ (meaning that if $U ⊂ V$ and $V$ satisfies $P$ then so does $U$). Assume
  the empty set satisfies $P$ and every point in $C$ has a neighborhood in $X$
  satisfying $P$. Then there exist sequences of subsets $K$ and $W$ indexed by
  natural numbers such that $K$ covers $C$, $W$ is locally finite and, for every $n$~:
  \begin{itemize}
    \item $K_n$ is compact
    \item $W_n$ is open
    \item $K_n ⊂ W_n$
    \item $W_n$ satisfies $P$.
  \end{itemize}
\end{lemma}

\begin{proof}
  \leanok
  This is a classical result.
\end{proof}

In the next lemma, $P₀$ is again a background local predicate satisfied by all
maps entering the discussion, and $P₁$ is the main target local predicate. We also use
an extra predicate $P₀'$ that enters the patching assumption in an asymmetric way
and will allows to deduce a relative version of the lemma.

\begin{lemma}
  \label{lem:inductive_construction_of_loc}\lean{inductive_construction_of_loc}\leanok
  Let $X$ a second countable locally compact metrizable topological space. Let $P₀$, $P₀'$
  and $P₁$ be local predicates on function from $X$ to a set $Y$.
  Let $f₀ \co X → Y$ be a function satisfying $P₀$ and $P₀'$. Assume that
  \begin{itemize}
    \item
      For every $x$ in $X$, there exists a function $f \co X → Y$ which
      satisfies $P₀$ and satisfies $P₁$ near $x$.
    \item
      For every closed subsets $K₁$ and $K₂$ of $X$ and every open subsets $U₁$ and $U₂$
      containing $K₁$ and $K₂$, for every function $f₁$ and $f₂$ satisfying $P₀$,
      if $f₁$ satisfies $P₀'$ and satisfies $P₁$ on $U₁$ and if $f₂$ satisfies $P₁$ on $U₂$
      then there exists $f$ which satisfies $P₀$ and $P₀'$, and satisfies $P₁$ near
      $K₁ ∪ K₂$ and coincides with $f₁$ near $K₁ ∪ U₂^c$.
  \end{itemize}
  Then there exists $f$ which satisfies $P₀$, $P₀'$ and $P₁$.
\end{lemma}

\begin{proof}
  \leanok\uses{lem:inductive_construction,lem:exists_locally_finite_subcover_of_locally,def:restrict_germ_predicate}
  The assumptions on the topology of $X$ and local existence of solutions
  allow to apply \Cref{lem:exists_locally_finite_subcover_of_locally} to get sequences
  of subsets $K$ and $U$ of $X$ indexed by natural numbers such that $K$ covers
  $X$, $U$ is locally finite and, for every $i$:
  \begin{itemize}
    \item $K_i$ is compact
    \item $U_i$ is open
    \item $K_i ⊂ U_i$
    \item there is a function $f \co X → Y$ which satisfies $P₀$ and satisfies $P₁$ on $U_i$.
  \end{itemize}

  We then apply \Cref{lem:inductive_construction} to the family of subsets $U$
  with local predicates $\hat{P}₀$ combining $P₀$ and $P₀'$ and $\hat{P}₁^i$ asking that
  $P₁$ holds near $K_i$, and the trivial family of auxiliary predicates $\hat{P}'₀$.

  For this we need to explain how the patching assumption of the current lemma
  implies the induction assumption of \Cref{lem:inductive_construction}. So we fix
  an index $i ∈ ℕ$ and a function $f$ which satisfies $\hat{P}₀$ and satisfies all
  $\hat{P}₁^j$ for $j < i$. We denote by $K$ the closed subset
  $\bigcup_{j < i} K_j$ and denote by $V$ an open neighborhood of $K$ such that
  $f$ satisfies $P₁$ on $V$. The patching assumption applied to $K$, $K_i$, $V$
  and $U_i$ with functions $f$ and the local solution on $U_i$ gives a suitable
  new function.
\end{proof}

From the above lemma we can deduce a version with only a base local predicate $P₀$ and
a target one $P₁$ and starting from a function which is already good near some closed
subset $K$. The is the version that is actually used in our application.

\begin{lemma}
  \label{lem:relative_inductive_construction_of_loc}
  \lean{relative_inductive_construction_of_loc}\leanok
  Let $X$ a second countable locally compact metrizable topological space. Let $P₀$
  and $P₁$ be local predicates on functions from $X$ to a set $Y$.
  Let $K$ be a closed subset of $X$.
  Let $f₀ \co X → Y$ be a function satisfying $P₀$ and satisfying $P₁$ near $K$.
  Assume that
  \begin{itemize}
    \item
      For every $x$ in $X$, there exists a function $f \co X → Y$ which
      satisfies $P₀$ and satisfies $P₁$ near $x$.
    \item
      For every closed subsets $K₁$ and $K₂$ of $X$ and every open subsets $U₁$ and $U₂$
      containing $K₁$ and $K₂$, for every function $f₁$ and $f₂$ satisfying $P₀$,
      if $f₁$ satisfies $P₁$ on $U₁$ and if $f₂$ satisfies $P₁$ on $U₁$
      then there exists $f$ which satisfies $P₀$, and satisfies $P₁$ near
      $K₁ ∪ K₂$ and coincides with $f₁$ near $K₁ ∪ U₂^c$.
  \end{itemize}
  Then there exists $f$ which satisfies $P₀$ and $P₁$ and coincides with $f₀$ near $K$.
\end{lemma}

\begin{proof}\uses{lem:inductive_construction_of_loc}\leanok
  We reduce this to \Cref{lem:inductive_construction_of_loc} using as auxilliary
  local predicate $P₀'$ the constraint to coincide with $f₀$ near $K$.

  Our patching condition almost matches the one from \Cref{lem:inductive_construction_of_loc}
  except that each $(K₁, K₂, U₁, U₂)$ should be replaced by $(K ∪ K₁, U ∪ U₁, K₂, U₂)$
  where $U$ is a suitable neighborhood of $K$.
\end{proof}

\begin{theorem}[Gromov]
  \label{thm:open_ample}
  % \uses{def:h-princ} % causes transitive arrow in depgraph
  \leanok
  \lean{RelMfld.Ample.satisfiesHPrincipleWith}
  For any manifolds $X$ and $Y$, any relation $\Rel ⊂ J^1(X, Y)$ that is open
  and ample satisfies the full $h$-principle (relative, parametric and $C^0$-dense).
\end{theorem}

\begin{proof}[Proof of Theorem~\ref{thm:open_ample}]
  \leanok
  \proves{thm:open_ample}
  \uses{lem:param_for_free, lem:ample_parameter, lem:transfer,
  lem:ex_localisation, lem:localisation_stability,
  lem:smooth_updating, lem:dist_updating, lem:inductive_htpy_construction,
  lem:ample_iff_loc, lem:h_principle_open_ample_loc}
  Lemmas~\ref{lem:param_for_free} and~\ref{lem:ample_parameter} prove we can
  assume there are no parameters. So we start with a single formal solution $F₀$
  of $\Rel$, which is holonomic near some closed subset $A ⊂ X$. We also
  fix a positive continuous function $δ$ on $X$ and we want to build a homotopy
  of formal solutions starting at $F₀$ relative to $A$ with base map staying at
  distance at most $δ$ from the base map of $F₀$ and ending at a holonomic
  solution.

  We apply \cref{lem:ex_localisation}  to get some
  localisation data $(φ \co \IT{N} × E → \set{X}, ψ \co ι × E' → \set{Y}, j)$ for
  $\bs F₀ : X → Y$. \Cref{lem:localisation_stability} then provides a continuous
  function $ε : X → ℝ_{> 0}$ such that every function $g$ with $d(\bs F₀, g) < ε$
  sends each $φ_i(E)$ into $ψ_{j(i)}(E')$. We denote by $τ$ the positive
  continuous function $\min(δ, ε)$.

  We then use the inductive construction of homotopies provided by
  \Cref{lem:inductive_htpy_construction} starting with $F₀$ and
  using the following local predicates. On maps $F$ from $X$ to $J¹(X, Y)$ we
  use the background predicate $P₀$ asserting that $F$ is a smooth section of
  $J¹(X, Y) → X$ that takes values in $R$, coincides with $F₀$ near $A$ and
  satisfies $d(\bs F, \bs F₀) < τ$. The background
  predicate for maps from $ℝ × X$ to $J¹(X, Y)$ is simply smoothness and the
  target local predicate $P₁$ on maps from $X$ to $J¹(X, Y)$ is being holonomic.
  We use the family of sets $U \co i ↦ φ_i(E)$ and $K \co i ↦ φ_i(\bar{B}_E)$.

  In order to check the induction assumption from
  \Cref{lem:inductive_htpy_construction}, we fix some
  $i$ in $\IT{N}$, and some formal solution $F$ which coincides with $F₀$ near
  $A$ and such that $d(\bs F₀, \bs F) < τ$. We assume that $F$ is holonomic near
  $\bigcup_{j < i} K_j$. We need to build a smooth homotopy of formal solutions
  starting at $F$ which coincide with $F₀$ near $A$, coincide with $F$ outside
  $U_i$, have base map at distance less than $τ$ from $\bs F₀$ and end at a
  formal solution which is holonomic near $\bigcup_{j ≤ i} K_j$. In addition this
  homotopy must be time-independent for time near $(-∞, 0]$ and $[1, +∞)$.

  Of course this homotopy comes from the local $h$-principle we proved in
  \Cref{lem:h_principle_open_ample_loc}. The first key observation allowing to
  apply that lemma is that $d(\bs F, \bs F₀) < τ ≤ ε$ hence $\bs F$ sends
  sends $φ_i(E)$ into $ψ_{j(i)}(E')$.

  \Cref{def:transfer_map} then turns $F$ into a section $\F$ of $J¹(E, E')$.
  According to \Cref{lem:transfer}, $\F$ is a formal solution of
  the relation $\Rel_i$ induced by $\Rel$ in $J¹(E, E')$ via $φ_i$ and
  $ψ_{j(i)}$, $\F$ is relative to $φ_i⁻¹(A)$
  and $\F$ is holonomic near $φ_i⁻¹(A ∪ \bigcup_{j < i} φ_j(\bar B_E))$.

  The homotopy $H$ will be constructed by updating $F$ using some
  homotopy $\HH$ of sections of $J¹(E, E')$ with support in the closed ball $2B_E$
  and time independent for $t$ near $(-∞, 0] ∪ [1, +∞)$ (here by support we mean
  the closure of the set of points where $\HH_t$ differs from $\F$).
  In order to ensure
  $d(\bs F₀, \bs H_t) < τ$ for all $t$, it suffices to ensure that, for each
  $x ∈ φ_i(2\bar{B}_E)$ and $t ∈ [0, 1]$,
  $d(\bs H_t(x), \bs F(x)) < τ(x) - d(\bs F(x), \bs F₀(x))$.
  The latter will hold as soon as, for all $e$ and $t$,
  $‖\bs \HH_t(e) - \bs \F(e)‖ < η$ for some positive $η$ given by
  \Cref{lem:dist_updating} (applied to $P = ℝ$, $M$ and $N$).
\end{proof}

\end{document}
